\documentclass[11pt,a4paper]{article}

\usepackage[T1]{fontenc}
\usepackage{amsmath,amssymb,amsthm,mathtools}
\usepackage[margin=25mm]{geometry}
\usepackage{lmodern,microtype}
\usepackage{booktabs,array,longtable}
\usepackage{enumitem}
\usepackage[hidelinks]{hyperref}
\usepackage{xurl}

\hypersetup{
  pdftitle={Finite-Fringe Rigidity in Congruence-Restricted Greedy Dissociated Sequences},
  pdfauthor={ProveIt research report (AI-assisted), merged from two manuscripts prepared for Vladimir Reshetnikov with OpenAI}
}

\setlength{\parskip}{3pt}
\setlength{\emergencystretch}{2em}
\setlength{\LTcapwidth}{\textwidth}
\numberwithin{equation}{section}

\newtheorem{theorem}{Theorem}[section]
\newtheorem{lemma}[theorem]{Lemma}
\newtheorem{proposition}[theorem]{Proposition}
\newtheorem{corollary}[theorem]{Corollary}
\theoremstyle{definition}
\newtheorem{definition}[theorem]{Definition}
\theoremstyle{remark}
\newtheorem{remark}[theorem]{Remark}
\theoremstyle{plain}
\newtheorem{question}{Question}

\newcommand{\Z}{\mathbb Z}
\newcommand{\N}{\mathbb N}
\newcommand{\cD}{\mathcal D}
\newcommand{\cE}{\mathcal E}
\newcommand{\cS}{\mathcal S}
\newcommand{\cF}{\mathcal F}
\newcommand{\ind}{\mathbf 1}
\newcommand{\res}{\operatorname{res}}
\newcommand{\Span}{\operatorname{span}}
\newcommand{\oeis}[1]{\href{https://oeis.org/#1}{OEIS #1}}
\newcommand{\mergenote}[1]{\par\smallskip\noindent\textbf{[Merge note, batch 73O1.]}\ #1\par\smallskip}
\newcommand{\file}[1]{\nolinkurl{#1}}
\newcommand{\src}[1]{\textup{(source #1)}}

\title{Finite-Fringe Rigidity in Congruence-Restricted\\
Greedy Dissociated Sequences\\[6pt]
\large Proofs, exact formulas, and inverse transseries for OEIS A139217 and A139218\\[4pt]
\normalsize Merged from two independent manuscripts (batch 73O1, sources 50 and 43)}
\author{ProveIt research report (AI-assisted)\\
prepared for Vladimir Reshetnikov with OpenAI}
\date{1 October 2026}

\begin{document}
\maketitle

\begin{abstract}
A finite set of positive integers is \emph{dissociated} when all of its subset
sums are distinct.  OEIS A139217 and A139218 are obtained by a greedy rule:
at each step choose the least admissible positive integer congruent to $1$ or
$2$ modulo $3$, respectively.  Both entries conjecture an eventual period-three
law for $a(n+1)-2a(n)$ and the same order-three linear recurrence.

We prove substantially stronger statements.  After three terms, the signed
subset-sum difference set is a complete interval with a fixed finite fringe of
holes.  For A139217 the only positive fringe hole is $S_n-3$; for A139218 the
positive fringe holes are $S_n-h$ with $h\in\{1,3,6,11\}$, where
$S_n=\sum_{j\le n}a(j)$.  A general finite-fringe propagation theorem shows
that these descriptions persist forever.  The greedy choice then becomes a
finite residue calculation, giving exact period-three corrections, rational
generating functions, closed forms, asymptotics, logarithmic transseries, and
residue-branch inverse transseries.

In particular, for A139217,
\[
 a(n)=\frac34\,2^n+
 \begin{cases}1,&n\equiv0,1\pmod3,\\-2,&n\equiv2\pmod3,
 \end{cases}
 \qquad(n\ge3),
\]
and for A139218,
\[
 a(n)=\frac{39}{56}\,2^n+
 \begin{cases}-25/7,&n\equiv0\pmod3,\\20/7,&n\equiv1\pmod3,\\5/7,&n\equiv2\pmod3,
 \end{cases}
 \qquad(n\ge4).
\]
Both satisfy $a(n)=a(n-1)+a(n-2)+2a(n-3)$ on their correct eventual ranges.
For A139217 that range is $n\ge6$; the current OEIS statement ``$n>4$'' is
false at $n=5$.  Self-contained exact verifiers are supplied.  The proofs are
conventional mathematics, not a formal Lean development, and independent
review is recommended before OEIS submission.

This report merges two manuscripts that proved these theorems independently
on the same day: source 50 (the base text, with the general finite-fringe
propagation theorem) and source 43 (boundary-defect coordinates).  Each shared
theorem is printed once and credited to both sources; source 43's different
proofs are kept as marked second routes, and its additional results are
printed in full: exact cross-identities between the two sequences, the sizes
of the signed difference sets, an explicit remainder bound and explicit
coefficients for the inverse series, counting asymptotics, and tables of
terms and defects for $n\le14$.
\end{abstract}

\clearpage
\setcounter{tocdepth}{1}
\tableofcontents
\clearpage

\section{Introduction}\label{gdf:sec:intro}

\subsection{The two OEIS problems}

Let $A_n=\{a(1),\dots,a(n)\}$.  The defining condition in the two entries is
that every two subsets of $A_n$ have different sums.  In modern additive
terminology, $A_n$ is a \emph{dissociated set}.  The greedy rule is constrained
to one arithmetic progression:
\begin{align*}
\text{A139217:}&\quad a(n+1)\text{ is the least positive }b\equiv1\pmod3
\text{ for which }A_n\cup\{b\}\text{ is dissociated},\\
\text{A139218:}&\quad a(n+1)\text{ is the least positive }b\equiv2\pmod3
\text{ for which }A_n\cup\{b\}\text{ is dissociated}.
\end{align*}
The entries currently list
\begin{align*}
\text{A139217: }&1,4,7,13,22,49,97,190,385,769,1534,3073,6145,12286,\dots,\\
\text{A139218: }&2,5,8,14,23,41,92,179,353,716,1427,2849,5708,11411,\dots.
\end{align*}
Both entries conjecture that the deviations $a(n+1)-2a(n)$ are eventually
periodic and that
\begin{equation}\label{gdf:eq:common-recurrence}
 a(n)=a(n-1)+a(n-2)+2a(n-3)
\end{equation}
eventually.  See \oeis{A139217} and \oeis{A139218}
\cite{OEIS217,OEIS218}, accessed 1 October 2026.
The entries were contributed by John W. Layman in April 2008.  Both also
mention the analogous odd-integer construction, which produces a shifted
Jacobsthal sequence.  As of the access date the order-three formulas are
labelled ``It appears'', with stated ranges $n>4$ and $n>6$ respectively
\src{43}.

The aim of this report is not merely to verify more terms.  We identify a
stable structural invariant of the complete signed difference set.  Once that
invariant is found, the conjectured recurrences become short consequences of
a finite residue table.

\subsection{Main results}

Write
\[
 S_n:=\sum_{j=1}^n a(j),\qquad
 \cD_n:=\left\{\sum_{j=1}^n\varepsilon_j a(j):
          \varepsilon_j\in\{-1,0,1\}\right\}.
\]
The strongest concrete conclusions are as follows.  Both theorems were proved
independently by sources 50 and 43 (source 43's Theorems~4.1 and~5.1).

\begin{theorem}[A139217, full signed support; sources 50 and 43]\label{gdf:thm:intro-217}
For every $n\ge3$,
\[
 \cD_n=[-S_n,S_n]\cap\Z\setminus\{-(S_n-3),S_n-3\}.
\]
Consequently,
\[
 a(n+1)-2a(n)=
 \begin{cases}
 -1,&n\equiv0\pmod3,\\
 -4,&n\equiv1\pmod3,\\
  5,&n\equiv2\pmod3,
 \end{cases}
 \qquad(n\ge3),
\]
and the recurrence \eqref{gdf:eq:common-recurrence} holds exactly for $n\ge6$.
\end{theorem}

\begin{theorem}[A139218, full signed support; sources 50 and 43]\label{gdf:thm:intro-218}
For every $n\ge3$,
\[
 \cD_n=[-S_n,S_n]\cap\Z\setminus
 \bigl\{\pm(S_n-h):h\in\{1,3,6,11\}\bigr\}.
\]
Consequently,
\[
 a(n+1)-2a(n)=
 \begin{cases}
 10,&n\equiv0\pmod3,\\
 -5,&n\equiv1\pmod3,\\
 -5,&n\equiv2\pmod3,
 \end{cases}
 \qquad(n\ge4),
\]
and the recurrence \eqref{gdf:eq:common-recurrence} holds exactly for $n\ge7$.
\end{theorem}

The full interval-minus-fringe formulas are stronger than the OEIS
conjectures: they determine every signed relation target, not only the next
greedy term.

Source 43 records the resolved OEIS statements as follows.
\begin{enumerate}[label=(\arabic*)]
\item For A139217, \eqref{gdf:eq:common-recurrence} holds for every $n\ge6$.
      The OEIS range $n>4$ must be corrected: at $n=5$,
      $22\ne13+7+2\cdot4=28$.
\item For A139218, \eqref{gdf:eq:common-recurrence} holds for every $n\ge7$,
      exactly matching the conjectured range $n>6$.
\item Both sequences have exact dyadic-plus-periodic closed forms, rational
      ordinary generating functions, and exact logarithmic staircase inverses.
\item The recurrence is a corollary of a stronger signed-difference-set
      classification with finitely many persistent boundary defects.
\end{enumerate}
The main proofs are elementary once the correct invariant is found.  The
nontrivial step is not solving a recurrence guessed from data, but proving that
the greedy construction continues to choose exactly the terms predicted by
that recurrence.  The signed difference set provides the missing certificate
\src{43}.

\subsection{Relation to recent work}\label{gdf:sec:recent}

Dissociated sets have a long history in distinct-subset-sum problems; useful
classical references include Cantor--Mills~\cite{CantorMills1966} and
Lunnon~\cite{Lunnon1988}, while modern additive terminology is discussed by
Candela--Helfgott~\cite{CandelaHelfgott2015} and Tao--Vu~\cite{TaoVu2006}.
The general distinct-subset-sums problem is a classical topic in additive
combinatorics; see, for example, Bohman's construction~\cite{Bohman1998}, the
lower-bound work of Dubroff--Fox--Xu~\cite{DubroffFoxXu2021}, and
Steinerberger's analytic reformulation~\cite{Steinerberger2023} \src{43}.
A recent paper of
Dutta~\cite{Dutta2026} studies the unrestricted greedy algorithm and proves
that, from any two fixed initial terms, the unrestricted greedy sequence
eventually doubles exactly.

The arithmetic-progression restriction changes the phenomenon.  Exact
doubling usually leaves the permitted residue class: if $a(n)\equiv1\pmod3$,
then $2a(n)\equiv2\pmod3$, and conversely.  The correct replacement is
\emph{doubling plus a bounded periodic correction}.  The finite-fringe theorem
below explains why; its threshold form (Corollary~\ref{gdf:cor:threshold})
is the explicit constrained analogue of Dutta's doubling \src{43}.

A targeted search on 1 October 2026 found the conjectures in the OEIS entries
and the general recent paper~\cite{Dutta2026}, but did not locate a published
proof of these two exact progression-restricted formulas.  This is a research
report, not a claim of priority after an exhaustive literature review.
Source 43's independent search used the exact OEIS identifiers, initial terms
and recurrence and likewise found no published proof; because such a search
can miss offline, unpublished or poorly indexed sources, its priority
statement is deliberately modest too.

\subsection{Status and reproducibility}\label{gdf:sec:status}

All arguments are written as ordinary proofs.  The accompanying
verifier \file{code/50-finite-fringe-verify.py} performs exact integer
checks of the greedy construction,
the fringe identities through 22 terms, the published prefixes, the closed
forms, the recurrences, and the rational generating functions.  The second
verifier \file{code/43-boundary-verify.py} (source 43) works with the signed
difference sets directly and also checks the counting functions, the
cross-identities and the failure at $n=5$
(Section~\ref{gdf:sec:audit}).  Computation is
used as an audit, not as a substitute for induction.  No theorem in this
report has been formalized in Lean at the time of writing.

Source 43 states its status as follows.  It is an unrefereed research draft.
The proofs are written to be independently checkable and are accompanied by
exact verification code, but they have not been formalized in Lean or
independently peer reviewed.  The literature and priority search was targeted
rather than exhaustive; no claim of global priority is made beyond the precise
statements and sources recorded here.  In particular source 43 does
\emph{not} claim:
\begin{itemize}
\item that no equivalent proof exists elsewhere;
\item that the general congruence-restricted greedy problem has been classified;
\item that the results have passed peer review or formal verification;
\item that exact computations alone establish the infinite claims.
\end{itemize}
The accompanying programs check finite prefixes and algebraic identities,
while the inductions in this report prove the statements for all indices.

The ProveIt repository was inspected at commit
\texttt{c79d64038558}~\cite{ProveIt} \src{50}.  The present topic is
separate from its existing A290268 development, but the report follows the
same useful pattern: an exact support model, a structural theorem, executable
checks, and an explicit formalization roadmap.
\mergenote{The A290268 development meant here is the collection report
\texttt{a290268-unbounded-deficits} of batch 72A; there is no mathematical
link between it and this report.}
The ProveIt repository combines formal mathematics, executable certificates,
and reproducible research artifacts, with a clear separation between
exploratory computation and trusted proof boundaries; the human-readable proof
is primary, and the verifiers are reproducibility aids \src{43}.  The inverse
expansions also fit the repository's broader series-and-transseries
programme~\cite{ProveItTransseries}, although the expansions obtained here are
unusually simple: on each residue branch they are convergent Taylor series in
$1/x$ \src{43}.

\subsection{The two sources and this merge}\label{gdf:sec:merge}

This report is built from two manuscripts of batch 73O1 of ProveIt's
incoming-reports intake, written independently on 1 October 2026.
\begin{itemize}
\item \textbf{Source 50} (the base text), \emph{Finite-Fringe Rigidity in
  Congruence-Restricted Greedy Dissociated Sequences}: the ternary support
  $\cE(A)$, the general finite-fringe propagation theorem with its collision
  condition, the residue rule for an arbitrary modulus, the periodic
  correction principle, logarithmic and ratio transseries, ternary seed
  certificates and ten research questions.  Its section order is kept.
\item \textbf{Source 43}, \emph{Boundary-Defect Rigidity in
  Congruence-Restricted Greedy Dissociated Sets}: the same fringe theorems,
  closed forms, recurrence ranges, generating functions, inverse branches and
  counting functions, proved in signed (boundary-defect) coordinates.
\end{itemize}
Every result, proof, example, remark, question and non-claim of both
manuscripts is printed.  A theorem proved by both is printed once, in source
50's wording, and credited ``sources 50 and 43''.  Where source 43's proof is
genuinely different, it follows as a ``Second proof (source 43)''.  Source
43's propagation theorem is a corollary of source 50's
(Corollary~\ref{gdf:cor:threshold}); its periodic forcing principle is
source 50's periodic correction principle restated for partial sums
(Proposition~\ref{gdf:prop:forcing}).  Material only source 43 has is marked
\src{43}: Section~\ref{gdf:sec:cross}, the forward asymptotics and spectral
form in Section~\ref{gdf:sec:unified}, the remainder bound and explicit
coefficients in Section~\ref{gdf:sec:inverse}, the counting asymptotics, the
second verifier, its proposed OEIS wording, its roadmap, its research
questions, and Appendix~\ref{gdf:app:tables}.  Text added in the merge is
marked \textbf{[Merge note, batch 73O1.]}.  Appendix~\ref{gdf:app:prov}
records the provenance and every choice the merge made.

\subsection{Notation across the two sources}\label{gdf:sec:notation}

The two manuscripts use different letters for the same notions, and two
letters with different meanings.  This report uses source 50's notation
throughout; source 43's material is translated as in the table.

\begin{center}
\small
\begin{tabular}{>{\raggedright\arraybackslash}p{4.0cm}>{\raggedright\arraybackslash}p{3.5cm}>{\raggedright\arraybackslash}p{2.1cm}>{\raggedright\arraybackslash}p{4.6cm}}
\toprule
Notion & This report & Source 50 & Source 43\\
\midrule
terms & $a(n)$; $u_n$, $v_n$ & $a(n)$ & $u_n$ (A139217), $v_n$ (A139218), $a_n$\\
partial sums & $S_n$; $U_n$, $V_n$ & $S_n$ & $U_n$, $V_n$, $S_n$\\
signed difference set & $\cD(A)$, $\cD_n$ & $\cD$ & $\Delta(A)$, $\Delta_n^{(u)}$, $\Delta_n^{(v)}$\\
ternary support & $\cE(A)$ & $\cE$ & ---\\
fringe offsets, $\max$ & $H$, $M=\max H$ & $H$, $M$ & $C$, $M$ \textbf{(renamed)}\\
interval with fringe & $\cF_H(S)$ & \eqref{gdf:eq:D-fringe-form} & $\mathcal F_C(S)$\\
greedy correction & $c_n=a(n+1)-S_n$ & $c_n$ & $-d_n$ (A139217), $-e_n$ (A139218) \textbf{(renamed, sign)}\\
doubling deviation & $d_n=a(n+1)-2a(n)$ & $d_n$ & no symbol\\
leading constant & $C$ & $C$ & $\alpha$ \textbf{(renamed)}\\
periodic part of $a(n)$ & $p_n$, $p_j$; $r_n$ for $v_n$ & $p_n$, $p_j$ & $p_n$ ($u$), $r_n$ ($v$), $p$\\
periodic part of $S_n$ & $q_n$ ($U$), $s_n$ ($V$) & --- & $q_n$, $s_n$\\
floor-log count & $\Lambda$ & $\Lambda$ & $\Phi$ \textbf{(renamed)}\\
inverse truncation order & $K$ & --- & $M$ \textbf{(renamed)}\\
forward shift operator & $\tau$ & --- & $E$ \textbf{(renamed)}\\
generating-function variable & $x$ & $x$ & $z$ \textbf{(renamed)}\\
\bottomrule
\end{tabular}
\end{center}

\noindent\textbf{Readings to avoid.}
\begin{itemize}
\item Source 43's $d_n$ is the \emph{offset} $S_n-a(n+1)$ of the greedy term
      below the partial sum, so $d_n^{(43)}=-c_n$.  It is \emph{not} source
      50's deviation $d_n=a(n+1)-2a(n)$.  For A139217, source 43's
      $d_n=(-1,3,-2)$ is $-c_n$ with $c_n=(1,-3,2)$, while the deviation is
      $d_n=(-1,-4,5)$.
\item Source 43's $C$ is a finite \emph{set} of offsets (this report's
      $H$); source 50's $C$, used here, is the leading \emph{constant} of
      $a(n)=C2^n+p_n$.
\item Source 43's $M$ is used both for $\max C$ and for a truncation order;
      here $M=\max H$ only, and the truncation order is $K$.
\item $\cE(A)$ is the ternary support; source 43's shift operator $E$ is
      written $\tau$.
\item The letter $p$ is source 50's period $p=m/\gcd(m,r)$
      (Section~\ref{gdf:sec:periodic}), while $p_n$ and $p_j$ denote periodic
      parts; a subscript always distinguishes them.
\item The word ``transseries'' in source 50's title and Section~\ref{gdf:sec:asymptotics}
      is used loosely: every expansion in this report is a convergent series
      in $2^{-n}$ or $1/x$, with no exponentially small sectors or
      trans-monomials.
\end{itemize}

\section{Dissociation, signed differences, and ternary supports}\label{gdf:sec:dissociation}

\subsection{Equivalent forms of dissociation}

For a finite set $A\subset\N$, define its ordinary subset-sum set and its
signed difference set by
\begin{align}
 \Sigma(A)&:=\left\{\sum_{a\in B}a:B\subseteq A\right\},\label{gdf:eq:sigma}\\
 \cD(A)&:=\left\{\sum_{a\in A}\varepsilon_a a:
                    \varepsilon_a\in\{-1,0,1\}\right\}.
 \label{gdf:eq:D-def}
\end{align}
Also write $S(A)=\sum_{a\in A}a$.

\begin{lemma}[standard equivalences]\label{gdf:lem:diss-equivalence}
For a finite set $A$ of positive integers, the following are equivalent.
\begin{enumerate}[label=(\roman*)]
\item All subset sums of $A$ are distinct.
\item The only relation $\sum_{a\in A}\varepsilon_a a=0$ with
      $\varepsilon_a\in\{-1,0,1\}$ is the trivial relation.
\item If $b>0$, then $A\cup\{b\}$ is dissociated exactly when
      $b\notin\cD(A)$, provided $A$ is dissociated.
\end{enumerate}
\end{lemma}

\begin{proof}
An equality between two subset sums becomes a signed relation after canceling
the common elements, and every signed relation can be separated into its
positive and negative supports.  This proves (i)$\Leftrightarrow$(ii).
For (iii), a relation involving the new element has coefficient $+1$ or $-1$
on $b$, hence is equivalent to $b\in\cD(A)$.  A relation with coefficient zero
on $b$ would already contradict dissociation of $A$.
\end{proof}

Statement (iii) is source 43's \emph{append criterion} (its Lemma~2.2), which
it calls decisive for the greedy algorithm.

\begin{proof}[Second proof of (iii) \src{43}]
The subset sums of $A\cup\{b\}$ split into two copies, $\Sigma(A)$ and
$b+\Sigma(A)$.  Each copy has no internal collision because $A$ is
dissociated.  A collision between the copies is an equality $b+s=t$ with
$s,t\in\Sigma(A)$, equivalently $b=t-s\in\Sigma(A)-\Sigma(A)=\cD(A)$.
Conversely, every representation $b=t-s$ creates that collision.
\end{proof}

The equality
\[
 \cD(A)=\Sigma(A)-\Sigma(A)
\]
is immediate.  In particular, $\cD(A)$ is symmetric and lies in
$[-S(A),S(A)]$.  Thus the greedy rule is not mysterious: at stage $n$, all
positive elements of $\cD(A_n)$ are forbidden, and the next term is the least
admissible integer outside that set \src{43}.

\begin{remark}[source 43]\label{gdf:rem:dutta-criterion}
Dutta uses the same signed-difference-set criterion in the unrestricted
greedy setting~\cite{Dutta2026}.  The new ingredient here is an exact
description of the whole support of $\cD(A_n)$ for the two
congruence-restricted OEIS sequences.
\end{remark}

\subsection{The shifted ternary support}

The central device is the support obtained by allowing digits $0,1,2$:
\begin{equation}\label{gdf:eq:E-def}
 \cE(A):=\left\{\sum_{a\in A}\delta_a a:
                \delta_a\in\{0,1,2\}\right\}.
\end{equation}

\begin{lemma}[translation identity]\label{gdf:lem:E-translation}
For every finite $A\subset\N$,
\[
 \cE(A)=S(A)+\cD(A).
\]
Consequently, $\cE(A)\subseteq[0,2S(A)]$, and $x\in\cE(A)$ if and only if
$2S(A)-x\in\cE(A)$.
\end{lemma}

\begin{proof}
Set $\delta_a=1+\varepsilon_a$.  Then
$\sum\delta_a a=S(A)+\sum\varepsilon_a a$.  Replacing every digit
$\delta_a$ by $2-\delta_a$ gives the reflection symmetry.
\end{proof}

Equivalently, $\cE(A)$ is the exponent support of
\begin{equation}\label{gdf:eq:support-polynomial}
 P_A(x):=\prod_{a\in A}(1+x^a+x^{2a}).
\end{equation}
Only the support matters here; coefficients may exceed one.

If $b$ is appended, the support has the exact three-copy recursion
\begin{equation}\label{gdf:eq:three-copy}
 \cE(A\cup\{b\})=\cE(A)\cup\bigl(b+\cE(A)\bigr)
                    \cup\bigl(2b+\cE(A)\bigr).
\end{equation}
This elementary identity drives the entire proof.  In signed coordinates it
reads $\cD(A\cup\{b\})=\cD(A)+\{-b,0,b\}$, the form source 43 uses.

\section{The finite-fringe propagation mechanism}\label{gdf:sec:propagation}

\subsection{Fringe-complete supports}

For a finite set $H\subset\N$ and an integer $T$, write
$T-H:=\{T-h:h\in H\}$ and $H+H:=\{h+h':h,h'\in H\}$.

\begin{definition}[finite-fringe completeness]\label{gdf:def:fringe}
Let $A$ have total $S$.  We say that $\cE(A)$ has fixed fringe $H$ if
\begin{equation}\label{gdf:eq:fringe-form}
 \cE(A)=([0,2S]\cap\Z)\setminus\bigl(H\cup(2S-H)\bigr).
\end{equation}
We always assume $H\subseteq\{1,\dots,S-1\}$.
\end{definition}

By Lemma~\ref{gdf:lem:E-translation}, \eqref{gdf:eq:fringe-form} is equivalent to
\begin{equation}\label{gdf:eq:D-fringe-form}
 \cD(A)=([-S,S]\cap\Z)\setminus
 \bigl\{\pm(S-h):h\in H\bigr\}.
\end{equation}
Thus the only positive admissible integers at most $S$ are the fringe gaps
$S-h$.  Source 43 names the right side of \eqref{gdf:eq:D-fringe-form} a
\emph{boundary-defect interval}: for $H$ finite and nonempty, $M=\max H$ and
$S>M$,
\begin{equation}\label{gdf:eq:FH}
 \cF_H(S):=([-S,S]\cap\Z)\setminus\{\pm(S-h):h\in H\},
\end{equation}
a complete symmetric integer interval with finitely many holes at fixed
distances $H$ from its two endpoints.

\subsection{The propagation theorem}

\begin{theorem}[finite-fringe propagation]\label{gdf:thm:fringe-propagation}
Let $A\subset\N$ have total $S$, and suppose \eqref{gdf:eq:fringe-form} holds for a
nonempty finite set $H$.  Put $M=\max H$.  Let $b$ be a positive integer such
that
\begin{equation}\label{gdf:eq:prop-conditions}
 M<b\le2S-M,
 \qquad
 2S-b\notin H+H.
\end{equation}
Then
\begin{equation}\label{gdf:eq:prop-conclusion}
 \cE(A\cup\{b\})=
 [0,2(S+b)]\cap\Z\setminus
 \bigl(H\cup(2(S+b)-H)\bigr).
\end{equation}
In words, appending $b$ preserves the same finite fringe.
\end{theorem}

\begin{proof}
Put $E=\cE(A)$ and consider the three translates in
\eqref{gdf:eq:three-copy}.  Their ambient intervals are
\[
 I_0=[0,2S],\qquad I_1=[b,b+2S],\qquad
 I_2=[2b,2b+2S].
\]
Because $b\le2S-M<2S$, these intervals cover the whole interval
$[0,2S+2b]$.

First consider a point covered only by $I_0$.  It is smaller than $b$.
The reflected holes $2S-H$ are at least $2S-M\ge b$, so the only holes that
can occur there are the desired left fringe $H$.  By reflection, a point
covered only by $I_2$ is missing exactly when it lies in the desired right
fringe $2(S+b)-H$.

A point covered only by $I_1$ can occur only when $b>S$.  Writing the point as
$b+y$, exclusivity gives
\[
 2S-b<y<b.
\]
The lower holes satisfy $h\le M\le2S-b$, while the upper holes satisfy
$2S-h\ge2S-M\ge b$.  Hence $y$ is not a hole of $E$, so every point in the
middle-only region is covered.

It remains to rule out an uncovered point in an overlap.  In an overlap of
two consecutive translates, such a point would require two holes of $E$ whose
difference is $b$.  Since $b>M$, two lower holes cannot differ by $b$, and two
upper holes cannot differ by $b$.  The only possible alignment is
\[
 2S-h=b+h'
 \quad(h,h'\in H),
\]
which is equivalent to $2S-b=h+h'\in H+H$.  This is excluded by
\eqref{gdf:eq:prop-conditions}.  Any overlap of $I_0$ and $I_2$ is automatically a
triple overlap, so a consecutive pair is available there as well.

Finally, $h\in H$ lies below $b$ and hence cannot be filled by a shifted copy;
similarly $2(S+b)-h$ lies beyond the end of $I_1$ and can only come from the
missing reflected point in $I_2$.  Therefore the endpoint holes remain, and
all other points are present.
\end{proof}

The collision condition has a transparent meaning: a high hole in one copy
must not land on a low hole in the next copy.

\begin{corollary}[automatic persistence beyond a threshold]
\label{gdf:cor:automatic-persistence}
Fix a modulus $m$ and a fringe $H$ with $M=\max H$.  Suppose a candidate has
the form $b=S+c$ with $-M\le c\le m$.  If
\[
 S>2M+m,
\]
then all conditions of Theorem~\ref{gdf:thm:fringe-propagation} hold.
\end{corollary}

\begin{proof}
We have $b\ge S-M>M$, and $b\le S+m<2S-M$.  Also
$2S-b=S-c\ge S-m>2M$, whereas $H+H\subseteq[2,2M]$.
\end{proof}

Thus a fixed fringe, once reached beyond a finite threshold, is rigid.  No
asymptotic estimate is needed.

Source 43 proves the threshold form of this persistence directly, in signed
coordinates and for an arbitrary correction.  It is the case of
Theorem~\ref{gdf:thm:fringe-propagation} recorded next (source 43's
Theorem~3.2, ``boundary-defect propagation'').

\begin{corollary}[threshold form; sources 50 and 43]\label{gdf:cor:threshold}
Let $H\subset\N$ be finite and nonempty, $M=\max H$, and let $c\in\Z$.
Suppose
\begin{equation}\label{gdf:eq:threshold}
 S>2M+|c|.
\end{equation}
Set
\[
 b=S+c,\qquad S'=S+b=2S+c.
\]
Then
\begin{equation}\label{gdf:eq:propagation-signed}
 \cF_H(S)+\{-b,0,b\}=\cF_H(S').
\end{equation}
Here $X+Y=\{x+y:x\in X,\ y\in Y\}$.  Consequently, if a dissociated prefix
$A_n$ with total $S_n$ satisfies $\cD(A_n)=\cF_H(S_n)$, the next chosen term is
$a(n+1)=S_n+c_n$, and $S_n>2M+|c_n|$, then
\begin{equation}\label{gdf:eq:inductive-certificate}
 \cD(A_{n+1})=\cF_H(S_{n+1}),
 \qquad S_{n+1}=2S_n+c_n
\end{equation}
(source 43's Corollary~3.3, ``inductive signed-difference certificate'').
\end{corollary}

\begin{proof}
By Lemma~\ref{gdf:lem:E-translation}, \eqref{gdf:eq:propagation-signed} is
\eqref{gdf:eq:prop-conclusion} shifted by $-(S+b)$, so it suffices to check
\eqref{gdf:eq:prop-conditions}.  We have $b\ge S-|c|>2M>M$ and
$b\le S+|c|<2S-2M\le 2S-M$.  Also $2S-b=S-c\ge S-|c|>2M\ge\max(H+H)$.  The
consequence follows because appending $a(n+1)$ changes the signed difference
set by $\cD(A_{n+1})=\cD(A_n)+\{-a(n+1),0,a(n+1)\}$.
\end{proof}

\begin{proof}[Second proof of \eqref{gdf:eq:propagation-signed} \src{43}]
First ignore the deleted points.  The three translated intervals are
\begin{align*}
[-S,S]-b&=[-S',-c],\\
[-S,S]&=[-S,S],\\
[-S,S]+b&=[c,S'].
\end{align*}
Because $S>|c|$, their union is all of $[-S',S']\cap\Z$.

We now track the holes.  Consider the right boundary point $S'-h$ with
$h\in H$.  It is larger than $S$ because $S+c-h>0$ follows from
\eqref{gdf:eq:threshold}.  Hence only the right translate can reach it.  Its
preimage there is $S-h$, which is deleted from $\cF_H(S)$.  Therefore $S'-h$
is absent.  By symmetry, $-S'+h$ is absent as well.

It remains to show that no other point is lost.  If $x\in[-S,S]$ is not a
central hole, the untranslated copy contains it.  If $x=S-h$ is a central
hole, then in the right translate it has preimage
\[
 x-b=(S-h)-(S+c)=-c-h.
\]
Now
\[
 |c+h|\le |c|+M<S-M.
\]
Every deleted point of $\cF_H(S)$ has absolute value at least $S-M$, so
$-c-h$ is present; hence $x$ is restored by the right translate.  The
negative central holes are restored symmetrically by the left translate.

Finally suppose $S<x\le S'$.  Only the right translate can reach $x$, and its
preimage $y=x-b$ lies in $(-c,S]$.  A negative hole has the form $-S+h$ and
satisfies
\[
 -S+h\le -S+M<-|c|\le -c,
\]
so a missing preimage in $(-c,S]$ can only be a positive hole $S-h$.  This
gives exactly $x=S'-h$, already listed.  The region $[-S',-S)$ is symmetric.
Therefore the only missing points are the prescribed boundary defects.
\end{proof}

\mergenote{Source 43 writes the new term as $a=S-d$ with hypothesis
$S>2M+|d|$; with $c=-d$ this is \eqref{gdf:eq:threshold}.  Its hypothesis
implies source 50's conditions \eqref{gdf:eq:prop-conditions}, but not
conversely: source 50's theorem also covers terms $b$ outside the window
$|b-S|<S-2M$.  The intake tested source 50's theorem in 11{,}984 random
fringe configurations and source 43's in 9{,}700, with no failure; source
43's hypothesis implied source 50's in every case.  In all 1{,}794 tested
cases with $M<b\le2S-M$ but $2S-b\in H+H$ the fringe broke, so the collision
condition cannot be dropped.}

\subsubsection*{Why this lemma is reusable \src{43}}

The theorem isolates a finite-state mechanism.  Once a greedy difference set
becomes a full interval with a fixed collection of boundary defects, the next
admissible term often has the form $S_n+c_n$, where $c_n$ depends only on a
residue state.  If $c_n$ is periodic, then $S_{n+1}=2S_n+c_n$ is exactly
solvable and the generated sequence is dyadic with a periodic correction.
Sections~\ref{gdf:sec:A139217} and~\ref{gdf:sec:A139218} are two concrete
realizations.

\subsection{Greedy choice from a fixed fringe}

Fix integers $m\ge2$ and $r\in\{0,1,\dots,m-1\}$.  Suppose every selected term
is congruent to $r$ modulo $m$.

\begin{proposition}[finite residue rule]\label{gdf:prop:residue-rule}
Assume \eqref{gdf:eq:fringe-form} with total $S>M$.  For each $h\in H$, the number
$S-h$ is admissible, and every other positive integer at most $S$ is
forbidden.  Therefore the next greedy term congruent to $r$ modulo $m$ is
\begin{equation}\label{gdf:eq:greedy-from-fringe}
 b=\min\left(
   \{S-h:h\in H,\ S-h\equiv r\pmod m\}
   \cup\{S+q\}\right),
\end{equation}
where $q\in\{1,\dots,m\}$ is the least positive integer satisfying
$S+q\equiv r\pmod m$.
\end{proposition}

\begin{proof}
Equation \eqref{gdf:eq:D-fringe-form} lists all admissible positive integers in
$[1,S]$.  Every integer above $S$ is outside $\cD(A)$ and hence admissible by
Lemma~\ref{gdf:lem:diss-equivalence}; $S+q$ is the first such integer in the
required residue class.
\end{proof}

If several $h$ have the required residue, the largest $h$ wins because it
gives the smallest $S-h$.  The correction
\begin{equation}\label{gdf:eq:c-def}
 c(S):=b-S
\end{equation}
is therefore selected from the finite set
$\{-h:h\in H\}\cup\{1,\dots,m\}$ and depends only on $S\bmod m$.

\subsection{Why a periodic correction gives a linear recurrence}\label{gdf:sec:periodic}

Let $S_n=\sum_{j\le n}a(j)$.  Since every term is $r$ modulo $m$,
\begin{equation}\label{gdf:eq:S-residue}
 S_n\equiv nr\pmod m.
\end{equation}
The residue cycle has period
\begin{equation}\label{gdf:eq:p-period}
 p=\frac{m}{\gcd(m,r)}.
\end{equation}
Under a permanent fixed fringe, Proposition~\ref{gdf:prop:residue-rule} therefore
gives a $p$-periodic sequence $c_n$ such that
\begin{equation}\label{gdf:eq:a-S-c}
 a(n+1)=S_n+c_n.
\end{equation}

\begin{theorem}[periodic correction principle]\label{gdf:thm:periodic-principle}
Suppose \eqref{gdf:eq:a-S-c} holds and $c_n$ is periodic with period $p$ from some
point onward.  Then
\begin{equation}\label{gdf:eq:d-c-difference}
 d_n:=a(n+1)-2a(n)=c_n-c_{n-1}
\end{equation}
is periodic with period $p$ and has zero sum over each period.  Consequently,
\begin{equation}\label{gdf:eq:p-recurrence}
 a(n+p)=a(n+p-1)+a(n+p-2)+\cdots+a(n+1)+2a(n)
\end{equation}
from the corresponding eventual range onward.  Its characteristic polynomial
is
\begin{equation}\label{gdf:eq:charpoly}
 (x-2)(1+x+\cdots+x^{p-1})
 =x^p-x^{p-1}-\cdots-x-2.
\end{equation}
\end{theorem}

\begin{proof}
Using $S_n=S_{n-1}+a(n)$ and
$a(n)=S_{n-1}+c_{n-1}$ gives
\[
 a(n+1)-2a(n)=S_n+c_n-2a(n)=c_n-c_{n-1}.
\]
The sum of $p$ consecutive first differences of the periodic sequence $c_n$
is zero.  Hence
\[
 0=\sum_{j=0}^{p-1}\bigl(a(n+j+1)-2a(n+j)\bigr),
\]
which simplifies to \eqref{gdf:eq:p-recurrence}.  Expanding the product gives
\eqref{gdf:eq:charpoly}.
\end{proof}

A further useful consequence is an exact exponential-plus-periodic form.
Since the roots other than $2$ in \eqref{gdf:eq:charpoly} are nontrivial
$p$th roots of unity, every solution is
\begin{equation}\label{gdf:eq:exp-periodic-general}
 a(n)=C2^n+p_n,
\end{equation}
where $p_n$ is $p$-periodic.  In the two OEIS cases $p=3$.

Source 43 proves the same principle in partial-sum coordinates, where the
greedy step \eqref{gdf:eq:a-S-c} reads $S_{n+1}=2S_n+c_n$.  Its statement
adds the uniqueness of the periodic part and the denominator of the
generating function.

\begin{proposition}[periodic forcing principle; sources 43 and 50]\label{gdf:prop:forcing}
Suppose that, from some index onward,
\begin{equation}\label{gdf:eq:periodic-forcing}
 S_{n+1}=2S_n+c_n,
 \qquad c_{n+p}=c_n
\end{equation}
for a fixed period $p$.  Then there is a unique $p$-periodic sequence $q_n$
satisfying
\[
 q_{n+1}=2q_n+c_n,
\]
and a constant $\kappa$ such that
\begin{equation}\label{gdf:eq:S-general-solution}
 S_n=\kappa 2^n+q_n.
\end{equation}
Consequently, $a(n)=S_n-S_{n-1}$ is a dyadic exponential plus a period-$p$
sequence of mean zero, in agreement with \eqref{gdf:eq:exp-periodic-general}
(with $C=\kappa/2$).  Its ordinary generating function is rational, with
denominator dividing
\begin{equation}\label{gdf:eq:general-denominator}
 (1-2x)(1+x+\cdots+x^{p-1}).
\end{equation}
\end{proposition}

\begin{proof}[Proof \src{43}]
Iterating the affine recurrence for $q$ over one period sends $q_n$ to
$2^pq_n+B_n$, where $B_n$ is determined by the $p$ forcing values.  A periodic
solution must satisfy $(1-2^p)q_n=B_n$, so it exists and is unique.
Subtracting this periodic solution from \eqref{gdf:eq:periodic-forcing} leaves
$S_{n+1}-q_{n+1}=2(S_n-q_n)$, proving \eqref{gdf:eq:S-general-solution}.

Now $a(n)=\kappa2^{n-1}+q_n-q_{n-1}$.  The periodic difference has mean zero
over one period, so its generating-function denominator divides
$(1-x^p)/(1-x)=1+x+\cdots+x^{p-1}$.  The dyadic part contributes $1-2x$.
\end{proof}

For $p=3$ the denominator is $(1-2x)(1+x+x^2)=1-x-x^2-2x^3$; equivalently the
characteristic polynomial is $(x-2)(x^2+x+1)=x^3-x^2-x-2$, which is exactly
the recurrence polynomial appearing in both OEIS entries.

\section{OEIS A139217: the residue-one case modulo 3}\label{gdf:sec:A139217}

\subsection{The seed fringe}

The first three greedy terms are $1,4,7$.  Indeed, after $1$ the values
$1$ and $-1$ are forbidden, so the next positive integer congruent to $1$
modulo $3$ is $4$; after $\{1,4\}$ the positive signed differences are
$1,3,4,5$, so the next permitted value is $7$.

Set $A_3=\{1,4,7\}$ and $S_3=12$.  From
\[
 (1+x+x^2)(1+x^4+x^8)(1+x^7+x^{14})
\]
one finds
\begin{equation}\label{gdf:eq:217-seed-E}
 \cE(A_3)=([0,24]\cap\Z)\setminus\{3,21\}.
\end{equation}
For completeness, the attainable values in $[0,12]$ are
\[
 0,1,2,4,5,6,7,8,9,10,11,12;
\]
reflection about $12$ supplies the upper half.  Thus the fringe is
\begin{equation}\label{gdf:eq:H217}
 H_{217}=\{3\}.
\end{equation}

\subsubsection*{Signed certificate \src{43}}

In signed coordinates \eqref{gdf:eq:217-seed-E} reads
\begin{equation}\label{gdf:eq:u-base}
 \cD(1,4,7)=[-12,12]\cap\Z\setminus\{\pm9\}=\cF_{\{3\}}(12).
\end{equation}
The positive integers $1,\dots,12$ other than $9$ have the signed
representations
\[
\begin{aligned}
 1&=1,        &2&=7-4-1, &3&=4-1,\\
 4&=4,        &5&=4+1,   &6&=7-1,\\
 7&=7,        &8&=7+1,   &10&=7+4-1,\\
 11&=7+4,     &12&=7+4+1.&&
\end{aligned}
\]
The value $9$ is impossible: with coefficient $0$ on $7$ the absolute value is
at most $5$; with coefficient $1$ on $7$, one would need a signed combination
of $1$ and $4$ equal to $2$, which does not occur; coefficient $-1$ cannot
produce a positive value as large as $9$.

\subsection{The residue table}

Every term is $1$ modulo $3$, so $S_n\equiv n\pmod3$.  The only positive gap
at most $S_n$ is $S_n-3$, which has the same residue as $S_n$.  Proposition
\ref{gdf:prop:residue-rule} gives the following exact table.

\begin{center}
{\small
\begin{tabular}{ccccc}
\toprule
$n\bmod3$ & $S_n\bmod3$ & next admissible $a(n+1)$ & $c_n=a(n+1)-S_n$ &
$d_n=a(n+1)-2a(n)$\\
\midrule
$0$ & $0$ & $S_n+1$ & $1$ & $-1$\\
$1$ & $1$ & $S_n-3$ & $-3$ & $-4$\\
$2$ & $2$ & $S_n+2$ & $2$ & $5$\\
\bottomrule
\end{tabular}
}
\end{center}

The $d_n$ column is $c_n-c_{n-1}$, with residues read cyclically.  Source 43
derives the same table: when $n\equiv1\pmod3$ the unique positive hole
$S_n-3$ is itself congruent to $1$ and is selected; in the other two cases the
hole has the wrong residue, every admissible positive integer at most $S_n$ is
represented, and the least admissible integer above $S_n$ is $S_n+1$ or
$S_n+2$.

\subsection{Persistence of the fringe}

Here $M=3$, $m=3$, and $S_3=12>2M+m=9$.  Every candidate in the table has
correction $c\in[-3,3]$.  Corollary~\ref{gdf:cor:automatic-persistence} applies
immediately.  Therefore the fringe \eqref{gdf:eq:H217} persists for every
$n\ge3$.  (In source 43's form: $12>2\cdot3+3$, so
Corollary~\ref{gdf:cor:threshold} applies from the base case onward for all
$c_n\in\{1,-3,2\}$.)

\begin{theorem}[complete solution of A139217; sources 50 and 43]\label{gdf:thm:A139217}
Let $a(n)$ be OEIS A139217 and $S_n=\sum_{j\le n}a(j)$.  Then:
\begin{enumerate}[label=(\alph*)]
\item For every $n\ge3$,
\[
 \cE(A_n)=[0,2S_n]\cap\Z\setminus\{3,2S_n-3\},
\]
and
\[
 \cD_n=[-S_n,S_n]\cap\Z\setminus\{\pm(S_n-3)\}.
\]
\item For every $n\ge3$,
\begin{equation}\label{gdf:eq:217-first-order}
 a(n+1)=2a(n)+
 \begin{cases}
 -1,&n\equiv0\pmod3,\\
 -4,&n\equiv1\pmod3,\\
 5,&n\equiv2\pmod3.
 \end{cases}
\end{equation}
\item For every $n\ge6$,
\begin{equation}\label{gdf:eq:217-recurrence}
 a(n)=a(n-1)+a(n-2)+2a(n-3).
\end{equation}
\item For every $n\ge3$,
\begin{equation}\label{gdf:eq:217-closed}
 a(n)=3\cdot2^{n-2}+
 \begin{cases}
 1,&n\equiv0,1\pmod3,\\
 -2,&n\equiv2\pmod3.
 \end{cases}
\end{equation}
Equivalently,
\begin{equation}\label{gdf:eq:217-cosine}
 a(n)=\frac34\,2^n-2\cos\left(\frac{2\pi(n+1)}3\right).
\end{equation}
\item The ordinary generating function $A_{217}(x)=\sum_{n\ge1}a(n)x^n$ is
\begin{equation}\label{gdf:eq:217-ogf}
 A_{217}(x)=
 \frac{x+3x^2+2x^3-6x^5}{1-x-x^2-2x^3}.
\end{equation}
\item For $n\ge3$,
\begin{equation}\label{gdf:eq:217-S}
 S_n=3\cdot2^{n-1}+
 \begin{cases}
 0,&n\equiv0\pmod3,\\
 1,&n\equiv1\pmod3,\\
 -1,&n\equiv2\pmod3.
 \end{cases}
\end{equation}
\end{enumerate}
\end{theorem}

\begin{proof}
Part (a) follows by induction from \eqref{gdf:eq:217-seed-E} and
Theorem~\ref{gdf:thm:fringe-propagation}, as justified above.  The residue table
then gives the three-periodic corrections $c_n=(1,-3,2)$ indexed by
$n\equiv(0,1,2)$.  Equation \eqref{gdf:eq:d-c-difference} gives (b).

The values $d_n$ in (b) sum to zero over a period.  Summing three consecutive
relations $a(k+1)-2a(k)=d_k$ gives (c), beginning with $n=6$ because the
periodic law already holds for $d_3,d_4,d_5$.

For (d), the claimed right side agrees with $a(3)=7$, and its periodic
correction $p_n\in\{1,1,-2\}$ satisfies
$p_{n+1}-2p_n=(-1,-4,5)$ in the corresponding residue classes.  Induction
using (b) proves the formula.  The cosine identity is the elementary
period-three interpolation of $(1,1,-2)$.

Multiplying $\sum_{n\ge1}a(n)x^n$ by
$1-x-x^2-2x^3$ and using (c) leaves only the coefficients through degree five,
which are $1,3,2,0,-6$; this gives (e).  Finally, summing (d), or checking the
three residue classes in $S_n=S_{n-1}+a(n)$, proves (f).
\end{proof}

In source 43's notation, part (a) with the greedy step reads: for every
$n\ge3$, $u_{n+1}=U_n+c_n$ with $c_n=(1,-3,2)$, hence
\begin{equation}\label{gdf:eq:U-first-order}
 U_{n+1}=2U_n+c_n,
\end{equation}
and part (f) is its Theorem~4.2, $U_n=3\cdot2^{n-1}+q_n$ with
$q_n=(0,1,-1)$ for $n\equiv(0,1,2)\pmod3$.

\begin{proof}[Second proof of (a), (d), (f) and (c) \src{43}]
Equation \eqref{gdf:eq:u-base} is the base case.  Assuming (a) at stage $n$,
the residue analysis proves that the greedy algorithm selects
$u_{n+1}=U_n+c_n$; the threshold condition holds, so
Corollary~\ref{gdf:cor:threshold} propagates the same defect set to stage
$n+1$.  This simultaneously proves the greedy choice and the invariant by
induction.

The periodic forcing in \eqref{gdf:eq:U-first-order} can be solved without
guessing.  A direct check of the three residue classes shows
$q_{n+1}=2q_n+c_n$.  Since $U_3=12=3\cdot2^2+q_3$, induction gives (f).
Taking $u_n=U_n-U_{n-1}$ gives the piecewise formula (d), including $n=3$ by
direct substitution; the trigonometric form is the real Fourier
representation of the period-three correction.

For (c), let $\tau$ be the forward shift.  The correction term $p_n$ is
annihilated by $\tau^2+\tau+1$, while $2^n$ is annihilated by $\tau-2$, so
$u_n$ is annihilated by $(\tau-2)(\tau^2+\tau+1)=\tau^3-\tau^2-\tau-2$.  The
closed form is valid from $n=3$, so applying its order-three annihilator
requires $n-3\ge3$, namely $n\ge6$.  At $n=5$ the two sides of
\eqref{gdf:eq:217-recurrence} are $22$ and $28$.
\end{proof}

Source 43 states (c) as its Corollary~4.3 and adds: ``This range is sharp for
the displayed initial sequence: the identity fails at $n=5$.''
\mergenote{The recurrence also holds at $n=4$ ($13=7+4+2\cdot1$), fails at
$n=5$, and holds for every $n\ge6$; $6$ is the first index from which it
holds permanently.  ``Sharp'' is true in that sense only.}

\subsection{Correction to the current OEIS range}

As of 1 October 2026, A139217 states that the recurrence holds for $n>4$.
At $n=5$ the proposed identity would read
\[
 22\stackrel{?}{=}13+7+2\cdot4=28,
\]
so that range is false.  Theorem~\ref{gdf:thm:A139217} shows that the correct
range is $n\ge6$.  The conjectured recurrence itself is correct after that
single indexing correction.  The eventual doubling correction recorded in the
OEIS entry is now exact: it is \eqref{gdf:eq:217-first-order}
(sources 50 and 43).

\section{OEIS A139218: the residue-two case modulo 3}\label{gdf:sec:A139218}

\subsection{The seed fringe}

The first three terms are $2,5,8$.  For $A_3=\{2,5,8\}$, $S_3=15$, and
\begin{equation}\label{gdf:eq:218-seed-E}
 \cE(A_3)=([0,30]\cap\Z)\setminus
 \{1,3,6,11,19,24,27,29\}.
\end{equation}
Indeed, the represented values in $[0,15]$ are
\[
 0,2,4,5,7,8,9,10,12,13,14,15,
\]
and reflection about $15$ gives the upper half.  Hence
\begin{equation}\label{gdf:eq:H218}
 H_{218}=\{1,3,6,11\}.
\end{equation}

\subsubsection*{Signed certificate \src{43}}

The positive represented signed values for $\{2,5,8\}$ are
\[
\begin{aligned}
 1&=8-5-2, &2&=2,       &3&=5-2,\\
 5&=5,     &6&=8-2,     &7&=5+2,\\
 8&=8,     &10&=8+2,    &11&=8+5-2,\\
 13&=8+5,  &15&=8+5+2.&&
\end{aligned}
\]
The four remaining positive integers are impossible.  Indeed,
$\cD(2,5)=\{0,\pm2,\pm3,\pm5,\pm7\}$; after fixing the coefficient of $8$ to
be $0$, $1$, or $-1$, none of $4,9,12,14$ can be completed by an element of
this set.  Hence
\begin{equation}\label{gdf:eq:v-base3}
 \cD(2,5,8)=[-15,15]\cap\Z\setminus
 \{\pm4,\pm9,\pm12,\pm14\},
\end{equation}
and $\{4,9,12,14\}=\{15-h:h\in H_{218}\}$.

\subsection{The residue table}

Now every term is $2$ modulo $3$, so
$S_n\equiv2n\pmod3$.  A gap $S_n-h$ is congruent to $2$ when
\[
 h\equiv S_n-2\equiv2(n-1)\pmod3.
\]
If several fringe elements match, the largest one yields the smallest greedy
candidate.  This gives:

\begin{center}
\begin{tabular}{cccccc}
\toprule
$n\bmod3$ & $S_n\bmod3$ & matching $h$ & chosen $h$ & $c_n$ & $d_n$\\
\midrule
$0$ & $0$ & $1$ & $1$ & $-1$ & $10$\\
$1$ & $2$ & $3,6$ & $6$ & $-6$ & $-5$\\
$2$ & $1$ & $11$ & $11$ & $-11$ & $-5$\\
\bottomrule
\end{tabular}
\end{center}

Thus the next term is always a fringe gap; no above-$S_n$ fallback is needed.
Source 43's table (``least admissible hole'' $v_{n+1}=V_n-1$, $V_n-6$,
$V_n-11$) is the same.

\subsection{Persistence of the fringe}

Here $M=11$.  At the initial stage $S_3=15$ and the table gives $b=14$.
The conditions of Theorem~\ref{gdf:thm:fringe-propagation} hold directly:
\[
 11<14\le19,
 \qquad
 2S_3-b=16\notin H_{218}+H_{218},
\]
where
\begin{equation}\label{gdf:eq:H218sum}
 H_{218}+H_{218}=\{2,4,6,7,9,12,14,17,22\}.
\end{equation}
After adjoining $14$, the new total is $S_4=29>2M+3=25$.
Corollary~\ref{gdf:cor:automatic-persistence} then applies forever.

\subsubsection*{The step $n=3\to4$ by hand \src{43}}

Source 43's threshold form does not cover the first step ($S_3=15<2\cdot11+1$),
so it computes $\cD(2,5,8,14)$ directly.  Appending $14$ replaces the signed
set by $\cD_3+\{-14,0,14\}$.  On the positive side, the translate by $14$
restores the four central holes because their preimages are
$-10,-5,-2,0\in\cD_3$.  Every integer from $16$ through $29$ has only the
preimage $x-14\in[2,15]$; it is missing precisely when that preimage is one of
$4,9,12,14$.  Symmetry now gives
\begin{equation}\label{gdf:eq:v-base4}
 \cD(2,5,8,14)
 =[-29,29]\cap\Z\setminus
 \{\pm18,\pm23,\pm26,\pm28\}
 =\cF_{H_{218}}(29).
\end{equation}
Thus both finite base identities have explicit hand-checkable certificates;
the accompanying program \file{code/43-boundary-verify.py} independently
verifies them by exact set equality.  For the induction step from $n=4$, the
threshold \eqref{gdf:eq:threshold} holds: $M=11$, $S_4=29$, and for the first
step $|c_4|=6$, so $29>22+6$; thereafter the partial sums increase, and the
largest possible $|c_n|$ is $11$.

\begin{theorem}[complete solution of A139218; sources 50 and 43]\label{gdf:thm:A139218}
Let $a(n)$ be OEIS A139218 and $S_n=\sum_{j\le n}a(j)$.  Then:
\begin{enumerate}[label=(\alph*)]
\item For every $n\ge3$,
\[
 \cE(A_n)=[0,2S_n]\cap\Z\setminus
 \bigl(H_{218}\cup(2S_n-H_{218})\bigr),
\]
where $H_{218}=\{1,3,6,11\}$.  Equivalently,
\[
 \cD_n=[-S_n,S_n]\cap\Z\setminus
 \bigl\{\pm(S_n-h):h\in H_{218}\bigr\}.
\]
\item For every $n\ge4$,
\begin{equation}\label{gdf:eq:218-first-order}
 a(n+1)=2a(n)+
 \begin{cases}
 10,&n\equiv0\pmod3,\\
 -5,&n\equiv1\pmod3,\\
 -5,&n\equiv2\pmod3.
 \end{cases}
\end{equation}
\item For every $n\ge7$,
\begin{equation}\label{gdf:eq:218-recurrence}
 a(n)=a(n-1)+a(n-2)+2a(n-3).
\end{equation}
\item For every $n\ge4$,
\begin{equation}\label{gdf:eq:218-closed}
 a(n)=\frac{39}{56}\,2^n+
 \begin{cases}
 -25/7,&n\equiv0\pmod3,\\
 20/7,&n\equiv1\pmod3,\\
 5/7,&n\equiv2\pmod3.
 \end{cases}
\end{equation}
Equivalently, in an explicitly integral piecewise form,
\begin{equation}\label{gdf:eq:218-closed-integral}
 a(n)=
 \begin{cases}
 (39\cdot2^n-200)/56,&n\equiv0\pmod3,\\
 (39\cdot2^n+160)/56,&n\equiv1\pmod3,\\
 (39\cdot2^n+40)/56,&n\equiv2\pmod3,
 \end{cases}
\end{equation}
that is, $a(n)=(39\cdot2^{n-3}-25)/7$, $(39\cdot2^{n-3}+20)/7$,
$(39\cdot2^{n-3}+5)/7$ respectively (source 43's form).
\item The ordinary generating function $A_{218}(x)=\sum_{n\ge1}a(n)x^n$ is
\begin{equation}\label{gdf:eq:218-ogf}
 A_{218}(x)=
 \frac{2x+3x^2+x^3-3x^4-9x^5-12x^6}
      {1-x-x^2-2x^3}.
\end{equation}
\item For every $n\ge3$,
\begin{equation}\label{gdf:eq:218-S}
 S_n=\frac{39}{28}\,2^n+
 \begin{cases}
 27/7,&n\equiv0\pmod3,\\
 47/7,&n\equiv1\pmod3,\\
 52/7,&n\equiv2\pmod3.
 \end{cases}
\end{equation}
\end{enumerate}
\end{theorem}

\begin{proof}
The direct base check and the propagation argument prove (a).  The residue
table gives $c_n=(-1,-6,-11)$ for $n\equiv(0,1,2)$, respectively.  Hence
$d_n=c_n-c_{n-1}=(10,-5,-5)$ once both adjacent corrections are in the stable
regime, namely for $n\ge4$.  This proves (b).

The three values of $d_n$ sum to zero, so Theorem~\ref{gdf:thm:periodic-principle}
gives (c), starting at $n=7$.  For (d), the periodic correction
$(-25/7,20/7,5/7)$ has successive doubled differences $(10,-5,-5)$, and the
formula agrees with $a(4)=14$; induction using (b) completes the proof.

Multiplication by $1-x-x^2-2x^3$ leaves the six initial residual coefficients
$2,3,1,-3,-9,-12$, giving (e).  Summation of (d), or induction in the three
residue classes, gives (f).
\end{proof}

In source 43's notation, part (a) with the greedy step reads: for every
$n\ge3$, $v_{n+1}=V_n+c_n$ with $c_n=(-1,-6,-11)$, hence
\begin{equation}\label{gdf:eq:V-first-order}
 V_{n+1}=2V_n+c_n,
\end{equation}
and part (f) is its Theorem~5.2, $V_n=\frac{39}{28}2^n+s_n$, with the
period-three sequence
\begin{equation}\label{gdf:eq:s-def}
 s_n=\begin{cases}
 27/7,&n\equiv0\pmod3,\\
 47/7,&n\equiv1\pmod3,\\
 52/7,&n\equiv2\pmod3,
 \end{cases}
\end{equation}
and $v_n=\frac{39}{56}2^n+r_n$, $r_n=(-25/7,20/7,5/7)$.

\begin{proof}[Second proof of (a), (d), (f) and (c) \src{43}]
The cases $n=3$ and $n=4$ are the direct computations
\eqref{gdf:eq:v-base3} and \eqref{gdf:eq:v-base4}; the transition $n=3$ to
$n=4$ also agrees with the greedy step.  For $n\ge4$, the residue table
identifies the greedy term, and Corollary~\ref{gdf:cor:threshold} propagates
the fixed defect set.  Induction completes the proof of (a).

One checks in the three residue classes that $s_{n+1}=2s_n+c_n$.  Since
$V_4=29=\frac{39}{28}2^4+\frac{47}7$, the first-order recurrence
\eqref{gdf:eq:V-first-order} gives (f) by induction.  Taking the difference
$V_n-V_{n-1}$ and simplifying gives (d); integrality also follows
automatically because the closed form equals the integer sequence generated
by the greedy construction.  For (c), the dyadic term and the period-three
correction are annihilated by $(\tau-2)(\tau^2+\tau+1)$; the closed form
begins at $n=4$, so the recurrence begins at $n=7$.
\end{proof}

\mergenote{Source 43 states (f) for $n\ge4$ only, starting the induction at
$V_4=29$.  The formula also holds at $n=3$: $\frac{39}{28}\cdot8+\frac{27}7=15=V_3$
(source 50 states $n\ge3$).  Source 43's delivered OEIS draft
(\file{43-boundary-OEIS_update_draft.txt}, shipped unchanged) keeps
``For $n>=4$''.}

The current OEIS range $n>6$ for \eqref{gdf:eq:218-recurrence} is therefore
correct.  The exact doubling correction \eqref{gdf:eq:218-first-order} proves
the eventual periodicity observed in the OEIS entry.

\section{Relations between the two sequences \src{43}}\label{gdf:sec:cross}

\subsection{Cross-sequence identities}

The dominant constants satisfy
\[
 \frac{39/56}{3/4}=\frac{13}{14}.
\]
Subtracting the two closed forms \eqref{gdf:eq:217-closed} and
\eqref{gdf:eq:218-closed} gives a stronger exact statement.

\begin{corollary}[cross-sequence identity]\label{gdf:cor:cross-relation}
For every $n\ge4$,
\begin{equation}\label{gdf:eq:cross-relation}
 14v_n-13u_n=
 \begin{cases}
 -63,&n\equiv0\pmod3,\\
 27,&n\equiv1\pmod3,\\
 36,&n\equiv2\pmod3.
 \end{cases}
\end{equation}
In particular,
\[
 \frac{v_n}{u_n}=\frac{13}{14}+O(2^{-n}).
\]
\end{corollary}

\begin{proof}
For $n\ge4$ both closed forms hold, and $14\cdot\frac{39}{56}=13\cdot\frac34$,
so the dyadic terms cancel.  The periodic remainders are
$14\cdot(-\frac{25}7)-13\cdot1=-63$,
$14\cdot\frac{20}7-13\cdot1=27$ and $14\cdot\frac57-13\cdot(-2)=36$.
\end{proof}

The range is sharp: at $n=3$, $14\cdot8-13\cdot7=21\ne-63$.  A corresponding
identity for partial sums is
\begin{equation}\label{gdf:eq:cross-sums}
14V_n-13U_n=
\begin{cases}
54,&n\equiv0\pmod3,\\
81,&n\equiv1\pmod3,\\
117,&n\equiv2\pmod3,
\end{cases}
\qquad(n\ge4),
\end{equation}
by the same subtraction applied to \eqref{gdf:eq:217-S} and
\eqref{gdf:eq:218-S}.
\mergenote{Since both partial-sum formulas hold from $n=3$,
\eqref{gdf:eq:cross-sums} also holds at $n=3$: $14\cdot15-13\cdot12=54$.}

\subsection{Near-saturation of the signed-difference interval}

The structural theorems immediately give the cardinalities
\begin{align}
 |\cD_n|&=2U_n-1\quad\text{(A139217)},\label{gdf:eq:u-delta-size}\\
 |\cD_n|&=2V_n-7\quad\text{(A139218)},\label{gdf:eq:v-delta-size}
\end{align}
for every $n\ge3$.  The ambient intervals have $2U_n+1$ and $2V_n+1$ points,
respectively.  Hence the support of the signed differences misses only $2$
points in the first sequence and $8$ in the second, independently of $n$.
This fixed-defect saturation is much stronger information than the term
recurrence.

\section{Exact asymptotics and transseries}\label{gdf:sec:asymptotics}

\subsection{Unified exponential-plus-periodic form}\label{gdf:sec:unified}

Both solutions have the exact form
\begin{equation}\label{gdf:eq:unified-form}
 a(n)=C2^n+p_j,\qquad j\equiv n\pmod3,
\end{equation}
with the following constants.

\begin{center}
\begin{tabular}{ccccc}
\toprule
Sequence & valid range & $C$ & $(p_0,p_1,p_2)$ & leading asymptotic\\
\midrule
A139217 & $n\ge3$ & $3/4$ & $(1,1,-2)$ & $a(n)\sim(3/4)2^n$\\
A139218 & $n\ge4$ & $39/56$ & $(-25/7,20/7,5/7)$ & $a(n)\sim(39/56)2^n$\\
\bottomrule
\end{tabular}
\end{center}

The periodic term is bounded, so there are no algebraic corrections in
$1/n$.  The complete correction scale is exponentially small relative to the
main term.

\begin{corollary}[forward asymptotics; source 43]\label{gdf:cor:forward-asymptotics}
As $n\to\infty$,
\begin{align*}
 u_n&=\frac34\,2^n+O(1),
 &U_n&=\frac32\,2^n+O(1),\\
 v_n&=\frac{39}{56}\,2^n+O(1),
 &V_n&=\frac{39}{28}\,2^n+O(1).
\end{align*}
The errors are exactly periodic of period $3$.  Moreover,
$u_{n+1}/u_n=2+O(2^{-n})$ and $v_{n+1}/v_n=2+O(2^{-n})$.
\end{corollary}

The closed forms are already complete asymptotic expansions: there are no
uncomputed algebraic or logarithmic correction terms.  They can also be viewed
as exact exponential polynomials \src{43}.  With $\omega=e^{2\pi i/3}$,
\[
 u_n=\frac34\,2^n-\omega^{n+1}-\omega^{-(n+1)},
\]
and the period-three term of $v_n$ is
\[
 r_n=-\frac{25}{7}\cos\frac{2\pi n}{3}
 +\frac{5\sqrt3}{7}\sin\frac{2\pi n}{3}.
\]
Thus the complete spectral support consists of the three characteristic roots
$2,\omega,\overline\omega$.

\subsection{Logarithmic transseries}

For a fixed residue branch $j$ and all sufficiently large $n\equiv j\pmod3$,
\begin{align}
 \log a(n)
 &=n\log2+\log C+
   \log\left(1+\frac{p_j}{C2^n}\right)\notag\\
 &=n\log2+\log C+
   \sum_{k\ge1}\frac{(-1)^{k+1}}{k}
   \left(\frac{p_j}{C}\right)^k2^{-kn}.
 \label{gdf:eq:log-transseries}
\end{align}
This is not merely asymptotic: it is a convergent expansion once
$|p_j|<C2^n$, which already holds throughout the closed-form ranges above.

Similarly, if $d_j$ is the period-three value of $a(n+1)-2a(n)$ on the branch
$n\equiv j$, then
\begin{align}
 \frac{a(n+1)}{a(n)}
 &=2+\frac{d_j}{C2^n+p_j}\notag\\
 &=2+\frac{d_j}{C}2^{-n}
   \sum_{k\ge0}\left(-\frac{p_j}{C}\right)^k2^{-kn}.
 \label{gdf:eq:ratio-transseries}
\end{align}
Thus convergence to the growth ratio $2$ has an explicit oscillatory
period-three expansion to all orders in $2^{-n}$.

\subsection{Inverse transseries on residue branches}\label{gdf:sec:inverse}

Because of the periodic correction, the inverse is naturally three-branched.
For a fixed residue $j$, solve $x=C2^n+p_j$ continuously for $n$:
\begin{equation}\label{gdf:eq:inverse-exact}
 \nu_j(x):=\log_2\left(\frac{x-p_j}{C}\right).
\end{equation}
Whenever $n$ lies in the closed-form range and $n\equiv j\pmod3$,
\begin{equation}\label{gdf:eq:inverse-exact-value}
 \nu_j(a(n))=n
\end{equation}
exactly.  Expanding at infinity gives the logarithmic inverse transseries
\begin{equation}\label{gdf:eq:inverse-transseries}
 \nu_j(x)=\log_2\left(\frac{x}{C}\right)
 -\frac1{\log2}\sum_{k\ge1}\frac{p_j^k}{k x^k}.
\end{equation}
For A139217 the leading inverse is $\log_2(4x/3)$; for A139218 it is
$\log_2(56x/39)$.  Formula \eqref{gdf:eq:inverse-transseries} supplies every
inverse-power correction and keeps the residue dependence explicit.  It was
obtained independently by both sources; for $|x|>|p_j|$ it is not merely a
formal Poincar\'e expansion but the convergent Taylor series of
$\log(1-p_j/x)$ \src{43}.

\begin{proposition}[truncation remainder; source 43]\label{gdf:prop:inverse-remainder}
Let $R_K(x)$ be the remainder after the first $K$ terms of the sum in
\eqref{gdf:eq:inverse-transseries}.  For $|x|>|p_j|$,
\begin{equation}\label{gdf:eq:inverse-remainder}
 \left|R_K(x)\right|
 \le
 \frac{|p_j|^{K+1}}{(K+1)\log2\,|x|^{K+1}}
 \frac{1}{1-|p_j/x|}.
\end{equation}
\end{proposition}

\begin{proof}
$|R_K(x)|\le\frac1{\log2}\sum_{k\ge K+1}\frac{|p_j/x|^k}{k}
\le\frac{1}{(K+1)\log2}\sum_{k\ge K+1}|p_j/x|^k$, a geometric series.
\end{proof}

\subsubsection*{Explicit branches \src{43}}

For A139217, $C=3/4$ and the periodic shifts are $p_j=1,1,-2$ for index
residues $j=0,1,2$.  Therefore
\begin{align}
 \nu_{0}(x)=\nu_{1}(x)
 &=\log_2\frac{4x}{3}
 -\frac1{\log2}\left(
 \frac1x+\frac1{2x^2}+\frac1{3x^3}+\frac1{4x^4}+\cdots\right),\label{gdf:eq:u-inverse-01}\\
 \nu_{2}(x)
 &=\log_2\frac{4x}{3}
 +\frac1{\log2}\left(
 \frac2x-\frac2{x^2}+\frac{8}{3x^3}-\frac4{x^4}+\cdots\right).
 \label{gdf:eq:u-inverse-2}
\end{align}
When $x=u_n$ and the appropriate branch is chosen, $\nu_j(x)=n$ exactly.

For A139218, $C=39/56$, with shifts $p_0=-\frac{25}{7}$, $p_1=\frac{20}{7}$,
$p_2=\frac{5}{7}$.  The first two corrections are
\begin{center}
\begin{tabular}{c c}
\toprule
$j$ & $\nu_j(x)-\log_2(56x/39)$\\
\midrule
$0$ & $\displaystyle \frac1{\log2}\left(\frac{25}{7x}-\frac{625}{98x^2}+O(x^{-3})\right)$\\[0.8em]
$1$ & $\displaystyle -\frac1{\log2}\left(\frac{20}{7x}+\frac{200}{49x^2}+O(x^{-3})\right)$\\[0.8em]
$2$ & $\displaystyle -\frac1{\log2}\left(\frac{5}{7x}+\frac{25}{98x^2}+O(x^{-3})\right)$\\
\bottomrule
\end{tabular}
\end{center}
Again, these are convergent for $|x|>|p_j|$.

\subsection{Exact counting functions}\label{gdf:sec:counting}

The closed forms also invert discretely.  Let
\[
 N_A(x):=\#\{n\ge1:a(n)\le x\}.
\]
Define, for real $y$,
\begin{equation}\label{gdf:eq:Lambda-def}
 \Lambda(y):=
 \begin{cases}
 \max\{0,\lfloor\log_8 y\rfloor\},&y>0,\\
 0,&y\le0.
 \end{cases}
\end{equation}
This counts positive integers $m$ for which $8^m\le y$.  (Source 43 calls this
function $\Phi$.)

For A139217, the three branches $n=3m,3m+1,3m+2$ give respectively
\[
 \frac{4(a(n)-1)}3=8^m,\qquad
 \frac{2(a(n)-1)}3=8^m,\qquad
 \frac{a(n)+2}3=8^m.
\]
Therefore, for every real $x$ (sources 50 and 43),
\begin{equation}\label{gdf:eq:count-217}
 \boxed{
 N_{217}(x)=\ind_{x\ge1}+\ind_{x\ge4}
 +\Lambda\!\left(\frac{4(x-1)}3\right)
 +\Lambda\!\left(\frac{2(x-1)}3\right)
 +\Lambda\!\left(\frac{x+2}3\right).}
\end{equation}

For A139218, the residue-zero branch begins with $n=6$ (that is, $m=2$),
while the other two closed-form branches begin with $m=1$.  Hence, for every
real $x$,
\begin{align}
 N_{218}(x)={}&\ind_{x\ge2}+\ind_{x\ge5}+\ind_{x\ge8}
 +\Lambda\!\left(\frac{7x+25}{39}\right)\notag\\
 &+\Lambda\!\left(\frac{28x-80}{39}\right)
 +\Lambda\!\left(\frac{14x-10}{39}\right).
 \label{gdf:eq:count-218}
\end{align}
Equations \eqref{gdf:eq:count-217}--\eqref{gdf:eq:count-218} are exact, not merely
asymptotic inversion formulas.

\mergenote{Source 50 printed the residue-zero term of
\eqref{gdf:eq:count-218} as
$\max\bigl\{0,\lfloor\log_8\bigl(\frac{56x+200}{39}\bigr)\rfloor-1\bigr\}$.
That expression equals $\Lambda\bigl(\frac{7x+25}{39}\bigr)$ for
$x>-25/7$, because $\frac{56x+200}{39}=8\cdot\frac{7x+25}{39}$, but it is
undefined for $x\le-25/7$, where the count is $0$.  The term printed here is
source 43's, which is defined for every real $x$.}

\begin{proof}[Proof \src{43}]
Besides the exceptional terms $u_1=1$ and $u_2=4$, \eqref{gdf:eq:217-closed}
gives, for $k\ge1$,
\[
 u_{3k}=\tfrac34\,8^k+1,\qquad
 u_{3k+1}=\tfrac32\,8^k+1,\qquad
 u_{3k+2}=3\,8^k-2.
\]
For a branch $A8^k+B$ with $k\ge1$, the number of terms not exceeding $x$ is
$\Lambda((x-B)/A)$.  Summing the three branches and the two exceptions gives
\eqref{gdf:eq:count-217}.  Besides $v_1=2,v_2=5,v_3=8$, the three branches of
A139218 may be written, for $k\ge1$, as
\[
 v_{3k+3}=\tfrac{39}{7}8^k-\tfrac{25}{7},\qquad
 v_{3k+1}=\tfrac{39}{28}8^k+\tfrac{20}{7},\qquad
 v_{3k+2}=\tfrac{39}{14}8^k+\tfrac{5}{7}.
\]
Applying the same branch count proves \eqref{gdf:eq:count-218}.
\end{proof}

\begin{corollary}[counting asymptotics; source 43]\label{gdf:cor:counting-asymptotics}
\[
 N_{217}(x)=\log_2 x+O(1),
 \qquad
 N_{218}(x)=\log_2 x+O(1)
 \qquad(x\to\infty).
\]
The $O(1)$ term cannot be replaced by a convergent constant: it contains
bounded logarithmically periodic staircase oscillations produced by the three
floor functions.
\end{corollary}

\subsection{Two notions of inversion \src{43}}

There are two complementary notions of inversion:
\begin{itemize}
\item the exact discrete counting functions \eqref{gdf:eq:count-217} and
      \eqref{gdf:eq:count-218}, which retain floors and bounded log-periodic
      oscillation;
\item the smooth residue-resolved inverses \eqref{gdf:eq:inverse-transseries},
      which admit all-orders convergent expansions.
\end{itemize}
Discarding the residue branch before inversion loses an $O(1)$ amount of
information, while discarding the periodic additive term within a fixed branch
loses only an explicitly expandable $O(x^{-1})$ correction to the index.

\section{Generating functions and recurrence diagnostics}\label{gdf:sec:gf}

The generating functions \eqref{gdf:eq:217-ogf} and \eqref{gdf:eq:218-ogf}
were found by both sources (source 43's Theorem~6.2, with variable $z$).
The common denominator factors as
\begin{equation}\label{gdf:eq:den-factor}
 1-x-x^2-2x^3=(1-2x)(1+x+x^2).
\end{equation}
The pole at $x=1/2$ produces the exponential term $C2^n$.  The roots of
$1+x+x^2$ are the nontrivial cube roots of unity and produce the bounded
period-three correction.  This analytic factorization exactly mirrors the
finite-fringe residue mechanism; the common denominator is not a numerical
coincidence (Proposition~\ref{gdf:prop:forcing}).

For A139217, the failure of the recurrence at $n=5$ is recorded by the term
$-6x^5$ in the numerator.  Beginning at degree six, multiplication by the
denominator gives zero.  For A139218, the transient lasts through degree six,
and the recurrence starts at degree seven.

The first-order deviations make the onset especially transparent:
\begin{center}
\begin{tabular}{ccl}
\toprule
Sequence & initial deviations $a(n+1)-2a(n)$ & periodic tail\\
\midrule
A139217 & $2,-1$ & $-1,-4,5,-1,-4,5,\dots$ from $n=3$\\
A139218 & $1,-2,-2$ & $-5,-5,10,-5,-5,10,\dots$ from $n=4$\\
\bottomrule
\end{tabular}
\end{center}
This reproduces and proves the patterns recorded as conjectural comments in
the OEIS entries.

\section{Exact computational audit}\label{gdf:sec:audit}

\subsection{Bitset model (source 50's verifier)}

The program \file{code/50-finite-fringe-verify.py} represents the support
of $P_A(x)$ in
\eqref{gdf:eq:support-polynomial} by one nonnegative Python integer.  Bit $j$ is
one exactly when $j\in\cE(A)$.  Appending $b$ is the exact operation
\begin{equation}\label{gdf:eq:bitset-update}
 B\longmapsto B\mathbin{\mathrm{OR}}(B\ll b)
                  \mathbin{\mathrm{OR}}(B\ll2b).
\end{equation}
No floating-point computation or probabilistic test is involved.

Given $B$ and $S$, positive admissible values $x\le S$ correspond exactly to
missing bits at positions $S+x$.  The program enumerates those missing bits,
filters by the desired congruence, and compares them with the first permitted
integer above $S$.  Thus it reconstructs the greedy sequence independently of
the closed forms.

\subsection{Checks performed by source 50's verifier}

Running
\begin{verbatim}
python code/50-finite-fringe-verify.py
\end{verbatim}
performs the following exact checks.
\begin{enumerate}[label=(\roman*)]
\item It generates 22 terms of each sequence directly from the greedy rule.
\item It matches the 14 terms currently printed by the OEIS.
\item From $n=3$ onward it verifies every missing exponent in
      $[0,2S_n]$, not merely a sample, against the claimed fixed fringe.
\item It checks the closed forms and recurrence on their stated ranges.
\item It expands both rational generating functions through degree 30 and
      checks every coefficient against the closed form.
\end{enumerate}
A representative run ends with
\begin{verbatim}
A139217: verified 22 greedy terms ... exact fringe [3].
A139218: verified 22 greedy terms ... exact fringe [1, 3, 6, 11].
A139217: rational generating function checked through x^30.
A139218: rational generating function checked through x^30.
All exact checks passed.
\end{verbatim}
It does not check the counting functions, the cross-identities or the failure
at $n=5$; source 43's verifier does.

\subsection{Source 43's verifier}

The program \file{code/43-boundary-verify.py} is a standard-library Python
program, intentionally independent of symbolic algebra systems.  All principal
checks use exact integers and rational numbers.  It performs the following
tasks.
\begin{enumerate}[label=(\arabic*)]
\item It generates the two sequences directly from the greedy definition,
      maintaining the exact signed difference set.
\item It verifies the published OEIS prefixes and extends each sequence to
      seventeen terms.
\item For every checked prefix, it proves by finite set equality that the
      signed difference set has exactly the boundary-defect form of
      Theorems~\ref{gdf:thm:A139217}(a) and~\ref{gdf:thm:A139218}(a).
\item It checks the closed forms and the correct recurrence ranges.
\item It explicitly confirms the failure of the A139217 recurrence at $n=5$.
\item It tests the boundary-propagation identity for all forcing values used
      in the proofs and several partial sums.
\item It checks the exact counting functions on every integer from $0$
      through the largest generated term.
\end{enumerate}
The core update is only
\begin{verbatim}
def signed_difference_set(values):
    out = {0}
    for value in values:
        out = {x + eps * value
               for x in out for eps in (-1, 0, 1)}
    return out
\end{verbatim}
The recorded output \file{data/43-boundary-verification_output.txt} begins:
\begin{verbatim}
All exact checks passed.
A139217 prefix: [1, 4, 7, 13, 22, 49, 97, 190, ...]
A139218 prefix: [2, 5, 8, 14, 23, 41, 92, 179, ...]
A139217 boundary offsets: {3}
A139218 boundary offsets: {1, 3, 6, 11}
\end{verbatim}

\subsection{Role and limits of computation}

The finite checks are useful for detecting transcription and indexing errors.
The proof of all $n$, however, is Theorem~\ref{gdf:thm:fringe-propagation} plus the
base fringes \eqref{gdf:eq:217-seed-E} and \eqref{gdf:eq:218-seed-E}.
Source 43 puts it the same way: the computation is useful for detecting the
invariant and guarding against transcription errors, but a finite calculation
cannot prove eventual behaviour.  The infinite argument consists of three
exact ingredients: a finite base difference-set identity; the greedy selection
forced by residues and boundary holes; and the propagation theorem.  Once
these are proved, no extrapolation from data remains.

\section{Proposed OEIS updates}\label{gdf:sec:oeis}

Both packages contain a plain-text draft suitable for adaptation after
independent review: \file{50-finite-fringe-oeis_update_draft.txt} and
\file{43-boundary-OEIS_update_draft.txt}.  Neither has been submitted.
They agree on every formula and range.  The mathematical content can be
summarized as follows.

\subsection{A139217}

\begin{itemize}
\item Replace the conjectural period statement by the proved identity
\[
 a(n+1)-2a(n)=(-1,-4,5)\quad\text{periodically for }n\ge3,
\]
with the values ordered by $n\equiv0,1,2\pmod3$.
\item Correct the recurrence range from ``$n>4$'' to ``$n\ge6$''.
\item Add the closed form \eqref{gdf:eq:217-closed}, generating function
\eqref{gdf:eq:217-ogf}, and exact signed-support statement from
Theorem~\ref{gdf:thm:A139217}.
\end{itemize}

\subsection{A139218}

\begin{itemize}
\item Replace the conjectural period statement by the proved identity
\[
 a(n+1)-2a(n)=(10,-5,-5)\quad\text{periodically for }n\ge4,
\]
with the values ordered by $n\equiv0,1,2\pmod3$.
\item Confirm the recurrence for $n\ge7$.
\item Add the closed form \eqref{gdf:eq:218-closed-integral}, generating function
\eqref{gdf:eq:218-ogf}, and exact signed-support statement from
Theorem~\ref{gdf:thm:A139218}.
\end{itemize}

The support theorems are compact enough for an OEIS comment but stronger than
what is needed to verify the numerical formula.  A stable public manuscript
link and checked attribution should be supplied before submission.

\subsection{Source 43's proposed wording}

Source 43 offers the following text as a starting point for an OEIS
submission, subject to independent review and attribution decisions.

\emph{A139217.} \textbf{Formula.} For $n\ge3$,
$a(n)=3\cdot2^{n-2}+1$ if $n\equiv0,1\pmod3$ and $a(n)=3\cdot2^{n-2}-2$ if
$n\equiv2\pmod3$.  Hence $a(n)=a(n-1)+a(n-2)+2a(n-3)$ for $n\ge6$.  The
previously displayed range $n>4$ is false at $n=5$.
\textbf{Comment.} If $S_n=\sum_{k=1}^n a(k)$, then for every $n\ge3$ the set
of signed sums
$\{\sum_{k=1}^n\varepsilon_k a(k):\varepsilon_k\in\{-1,0,1\}\}$ is
$[-S_n,S_n]\cap\Z$ with only $\pm(S_n-3)$ removed.  This invariant proves the
greedy formula.
\textbf{Generating function.} \eqref{gdf:eq:217-ogf}.

\emph{A139218.} \textbf{Formula.} For $n\ge4$, \eqref{gdf:eq:218-closed}.
Hence $a(n)=a(n-1)+a(n-2)+2a(n-3)$ for $n\ge7$.
\textbf{Comment.} If $S_n=\sum_{k=1}^n a(k)$, then for every $n\ge3$ the
signed-sum set is $[-S_n,S_n]\cap\Z$ with precisely
$\pm(S_n-1)$, $\pm(S_n-3)$, $\pm(S_n-6)$, $\pm(S_n-11)$ removed.
\textbf{Generating function.} \eqref{gdf:eq:218-ogf}.

Source 43's delivered draft file also gives the partial sums
\eqref{gdf:eq:217-S} and $S_n=\frac{39}{28}2^n+s_n$ (stated there for
$n\ge4$; true from $n=3$), source 43's integral form of
\eqref{gdf:eq:218-closed}, and a status and attribution note.  Source 50's
draft adds the cosine form \eqref{gdf:eq:217-cosine}, the periodic deviations,
the leading asymptotics and a proof idea.

\section{Formalization roadmap for ProveIt}\label{gdf:sec:roadmap}

The proof is well suited to Lean because its only infinite step is an
induction over finite integer intervals and finite set translations.

\subsection{Suggested definitions}

A direct formalization can use finite sets rather than polynomials:
\begin{align*}
 \texttt{signedSums}(A)&:=
 \left\{\sum_i\varepsilon_i a_i:\varepsilon_i\in\{-1,0,1\}\right\},\\
 \texttt{ternarySums}(A)&:=
 \left\{\sum_i\delta_i a_i:\delta_i\in\{0,1,2\}\right\}.
\end{align*}
For a list or finset of naturals, the second set can be defined recursively by
\[
 E([])=\{0\},\qquad E(a::L)=E(L)\cup(a+E(L))\cup(2a+E(L)).
\]
This recursion avoids coefficient reasoning entirely.

\subsection{Core lemmas}

A compact development would prove, in order:
\begin{enumerate}[label=(\arabic*)]
\item distinct subset sums $\Leftrightarrow$ no nontrivial signed zero
relation;
\item admissibility of a new term $b\Leftrightarrow b\notin\cD(A)$;
\item $\cE(A)=S(A)+\cD(A)$ and reflection symmetry;
\item the three-copy recursion \eqref{gdf:eq:three-copy};
\item Theorem~\ref{gdf:thm:fringe-propagation} as a lemma about three translated
integer intervals with finite deleted fringes;
\item the finite residue rule and periodic correction principle;
\item decidable base computations for $\{1,4,7\}$ and $\{2,5,8\}$;
\item the two inductions, closed forms, and recurrence identities.
\end{enumerate}

The only delicate part is organizing the interval-overlap proof so that
arithmetic automation sees the cases.  One practical approach is to state a
separate coverage lemma for $I_0,I_1,I_2$ and then discharge the possible
hole alignments with finite-set membership in $H+H$.  For the concrete
fringes, all remaining arithmetic should be handled by
\texttt{omega}, \texttt{norm\_num}, and finite \texttt{decide} checks.

\subsection{Proof-certificate boundary}

The executable verifier is intentionally not presented as a formal proof.
A Lean development should import neither the generated terms nor a large
bitset certificate as an axiom.  The finite seed supports may be proved by
normalization or reflected interval membership, while the universal step must
be the abstract fringe theorem.  This keeps the trusted kernel small and the
mathematics reusable for other OEIS sequences.

\subsection{Source 43's roadmap}

Source 43 proposes a compact Lean development with the following layers.

\emph{Core definitions.} Represent a finite prefix as a list or finset of
integers and define $\operatorname{delta}(A)=\{\sum_{a\in A}\varepsilon_a
a:\varepsilon_a\in\{-1,0,1\}\}$.  For formal efficiency, one can instead
define $\Sigma(A)$ first and set $\cD(A)=\Sigma(A)-\Sigma(A)$; the latter
makes the append criterion a direct statement about two translated finsets.
Suggested declarations include
\begin{verbatim}
subsetSums : List Int -> Finset Int
signedDiffs : List Int -> Finset Int
dissociated : List Int -> Prop
boundaryModel : Finset Nat -> Int -> Finset Int
\end{verbatim}

\emph{General theorem layer.} \nolinkurl{dissociated_cons_iff_not_mem_signedDiffs},
formalizing Lemma~\ref{gdf:lem:diss-equivalence}(iii); a translate identity
for signed differences after appending one term;
\texttt{boundaryModel\_propagate}, formalizing
Corollary~\ref{gdf:cor:threshold}; and interval-membership lemmas that
separate the central and two outer regions.  The propagation proof uses only
linear integer arithmetic, finite-set extensionality, and elementary
inequalities, so most local obligations should be amenable to
\texttt{omega}, \texttt{linarith}, and simplification.

\emph{Sequence-specific layer.} For A139217, certify the base equality for
$(1,4,7)$ and then prove the three residue cases.  For A139218, certify the
two finite bases $(2,5,8)$ and $(2,5,8,14)$.  These bases are small enough for
a transparent theorem-by-theorem proof; a checked decision procedure may also
be used if the repository's trust policy permits it.  After the structural
induction, the formulas are routine recurrences over residue classes.  The
final recurrence follows from either direct substitution in the piecewise
closed forms or an abstract lemma for dyadic-plus-periodic sequences.

\emph{Recommended theorem boundary.} A strong formalization should make the
main public theorems state the greedy semantics, not only the numerical
recurrence.  Proving that a recurrence-generated list matches the first
several terms is insufficient; the key statement is that every next term is
the least admissible integer in the required congruence class.  The
signed-difference invariant provides exactly that bridge.

\mergenote{None of these declarations exists in ProveIt.  The names above are
the sources' proposals; nothing in this report is formalized.}

\section{Further research questions}\label{gdf:sec:questions}

The fixed-fringe mechanism suggests a broader program rather than two isolated
identities.  Questions~\ref{gdf:q:classification}--\ref{gdf:q:oeis} are source
50's; Section~\ref{gdf:sec:questions43} lists source 43's, with the overlaps
marked.

\begin{question}[classification by modulus and residue]\label{gdf:q:classification}
For which pairs $(m,r)$ does the progression-restricted greedy dissociated
sequence eventually acquire a fixed fringe?  Can the fringe and the onset
index be computed from $(m,r)$ alone?
\end{question}

The cases $r=0$ reduce after scaling by $m$ to an unrestricted problem, while
nonzero residues can forbid exact doubling.  The first unexplored step is a
systematic census for $m=4,5,6$ with proof-oriented support data.

\begin{question}[a finite fringe automaton]\label{gdf:q:automaton}
Can one define a finite-state transition system whose states are normalized
hole sets near $0$ and $2S$, and prove that every progression-restricted
greedy orbit eventually enters a cycle?
\end{question}

A naive state space is infinite because a hole can drift into the bulk.  The
central issue is to prove a uniform bound on relevant fringe width, or to find
a counterexample with a growing fringe.  The three-copy recursion gives a
natural transfer operator on hole patterns.

\begin{question}[necessary and sufficient propagation conditions]\label{gdf:q:conditions}
Theorem~\ref{gdf:thm:fringe-propagation} gives simple sufficient conditions.
Can one classify exactly the values of $b$ for which a given fringe $H$ is
preserved?  What new fringe is produced when the collision
$2S-b\in H+H$ occurs?
\end{question}

The collision equation suggests a controlled ``fringe surgery'' calculus:
aligned high and low holes can either persist, be filled by the third copy, or
create a new finite cluster.  An exact transition theorem could make the
automaton question tractable.

\begin{question}[eventual exponential-plus-periodic laws]\label{gdf:q:laws}
Is every fixed-fringe progression sequence eventually of the form
$C2^n+p_n$ with periodic $p_n$?  More generally, if the normalized fringe is
periodic rather than fixed, does one still obtain a rational generating
function?
\end{question}

For a genuinely periodic family of fringes, the correction $c_n$ would remain
periodic after augmenting the state by the fringe phase.  This strongly
suggests rationality, but the onset and state bound need proof.

\begin{question}[higher representation multiplicity]\label{gdf:q:multiplicity}
Replace dissociation by the condition that every subset sum has at most $g$
representations, or replace digits $\{-1,0,1\}$ by
$\{-q,\dots,q\}$.  Is there an analogue of finite-fringe propagation with
$2q+1$ translated copies?
\end{question}

The geometric part of the proof generalizes immediately to multiple copies,
but overlap multiplicities and admissibility are subtler.  This direction
connects directly with the generalized dissociation questions considered in
\cite{Dutta2026}.

\begin{question}[unions of residue classes]\label{gdf:q:unions}
What happens when candidates are allowed in a union $R\subset\Z/m\Z$ rather
than one residue class?  Can the greedy choice be described by a finite
minimum over fringe gaps, and when does exact doubling reappear?
\end{question}

This interpolates between the present constrained sequences and Dutta's
unrestricted greedy algorithm.  It may reveal a sharp threshold between
exact eventual doubling and periodic near-doubling.

\begin{question}[robustness under seeds]\label{gdf:q:seeds}
Fix $(m,r)$ but begin with an arbitrary finite dissociated seed in the residue
class.  Which seeds converge to the same fringe cycle, and which produce
different asymptotic constants $C$?
\end{question}

The constant $C$ is determined by the transient together with the periodic
correction.  Basin-of-attraction questions may be finite-state once a fringe
bound is available.

\begin{question}[effective inverse theory]\label{gdf:q:inverse}
For eventual laws $a(n)=C2^n+p_n$, can one give uniform exact formulas for the
counting function, nearest-index map, and error bounds when only a recurrence
or fringe automaton is known?
\end{question}

The formulas in Section~\ref{gdf:sec:asymptotics} show that the inverse problem
naturally decomposes
by residue branch.  For longer periods, roots-of-unity filters and floor-log
expressions should give a systematic answer.

\begin{question}[formal library theorem]\label{gdf:q:formal}
Can Theorem~\ref{gdf:thm:fringe-propagation} be formalized once as a reusable
ProveIt lemma and applied automatically to a certificate consisting only of a
seed fringe and a finite residue table?
\end{question}

Such a theorem would turn many OEIS conjectures into small, auditable proof
certificates rather than sequence-specific formal developments.

\begin{question}[OEIS exploration]\label{gdf:q:oeis}
Which existing OEIS entries record unexplained eventual recurrences arising
from greedy distinct-subset-sum constructions?  Can their recurrences be
reverse-engineered into finite fringe sets?
\end{question}

A useful search strategy is to inspect entries with comments of the form
``$a(n+1)-2a(n)$ appears periodic,'' compute exact ternary-support holes, and
test whether the holes stabilize after translation by the total sum.

\subsection{Source 43's questions}\label{gdf:sec:questions43}

Source 43 numbers its questions consecutively; its numbers are kept.  For a
modulus $m\ge2$ and residue $r$, the greedy dissociated sequence chooses the
least positive integer congruent to $r$ modulo $m$ that preserves distinct
subset sums.

\emph{Classification by modulus and residue} (overlaps
Question~\ref{gdf:q:classification}).
\begin{enumerate}[label=43.\arabic*]
\item For which pairs $(m,r)$ does the signed difference set eventually become
      a full interval with a fixed finite boundary-defect set?
\item When this happens, can the defect set be computed from a finite seed by
      a terminating algorithm with proof certificates?
\item Is eventual boundary-defect stabilization guaranteed whenever
      $\gcd(m,r)=1$, or are there primitive counterexamples?
\end{enumerate}
Preliminary exact experiments indicate that fixed defect sets occur for
several further small pairs, but a systematic theorem requires more than
numerical evidence.

\emph{A finite-state classification theorem} (overlaps
Questions~\ref{gdf:q:automaton}--\ref{gdf:q:laws}).  The present examples
have a particularly rigid form: $\cD_n=\cF_H(S_n)$ and
$S_{n+1}=2S_n+c(S_n\bmod m)$.
\begin{enumerate}[label=43.\arabic*,resume]
\item Classify all finite sets $H$ and residue-choice functions $c$ for which
      the propagation equation \eqref{gdf:eq:propagation-signed} is
      compatible with a greedy rule.
\item Develop a finite automaton whose states encode the boundary defects and
      whose accepted cycles certify eventual recurrences.
\item Determine whether every ultimately periodic defect automaton yields a
      rational generating function, and characterize its minimal denominator.
\end{enumerate}

\emph{Beyond a single residue class} (overlaps
Question~\ref{gdf:q:unions}).
\begin{enumerate}[label=43.\arabic*,resume]
\item Replace one residue class by a union of classes, or by an arbitrary
      periodic admissibility set.  Does the next-step forcing become
      automatic, morphic, or eventually periodic?
\item Study admissible sets defined by valuations, digital constraints, or
      Beatty sequences.  Which constraints preserve finite-state boundary
      dynamics?
\item Allow a weighted greedy objective, such as the least admissible term
      under a cost function rather than the usual order.
\end{enumerate}

\emph{Robustness and seed dependence} (overlaps
Question~\ref{gdf:q:seeds}).
\begin{enumerate}[label=43.\arabic*,resume]
\item If finitely many initial terms are changed while preserving dissociation
      and congruence, does the process converge to the same defect cycle, a
      translated cycle, or one of finitely many attractors?
\item Quantify a basin of attraction for each stable defect set.
\item Determine the earliest stabilization index and sharp versions of the
      sufficient threshold $S>2M+|c|$ of Corollary~\ref{gdf:cor:threshold}.
\end{enumerate}
\mergenote{Item 43.12 is partly answered by source 50: its
Theorem~\ref{gdf:thm:fringe-propagation} replaces the threshold by the
conditions $M<b\le2S-M$, $2S-b\notin H+H$, whose collision condition is
sharp in the intake's random tests (see the note after
Corollary~\ref{gdf:cor:threshold}).  The
earliest stabilization index for general $(m,r)$ remains open.}

\emph{Connections to broader dissociated-set theory} (item 43.14 overlaps
Question~\ref{gdf:q:multiplicity}).
\begin{enumerate}[label=43.\arabic*,resume]
\item Dutta proved eventual exact doubling for the unrestricted greedy
      algorithm~\cite{Dutta2026}.  Find a common theorem that treats
      unrestricted doubling and congruence-restricted periodic corrections as
      two instances of one dynamical principle.
\item Extend the constrained greedy analysis to $D[g]$ sets, where each subset
      sum may have at most $g$ representations, and to $D_k$ sets with
      coefficients in $\{-k,\ldots,k\}$.
\item Compare the asymptotic constants produced by residue restrictions with
      extremal constructions in the Erd\H{o}s distinct-subset-sums
      problem~\cite{Bohman1998,DubroffFoxXu2021,Steinerberger2023}.
\item Determine exact representation multiplicities inside the nearly
      saturated difference intervals.  The support is simple, but the number
      of signed representations of each interior point may contain additional
      structure.
\end{enumerate}

\emph{Asymptotic inversion and transseries} (overlaps
Question~\ref{gdf:q:inverse}).
\begin{enumerate}[label=43.\arabic*,resume]
\item Develop a general inversion calculus for sequences of the form
      $\alpha\lambda^n+p(n)$ with periodic or automatic $p(n)$, retaining both
      exact staircase effects and smooth branch expansions.
\item For more complicated defect automata, identify when inverse expansions
      involve several exponential scales, logarithms, or genuine transseries
      rather than convergent $1/x$ series.
\item Formalize a theorem that transports error bounds from a forward
      exponential-polynomial approximation to residue-resolved inverse
      branches.
\end{enumerate}

\emph{Formal verification and OEIS infrastructure} (overlaps
Questions~\ref{gdf:q:formal} and~\ref{gdf:q:oeis}).
\begin{enumerate}[label=43.\arabic*,resume]
\item Formalize the propagation theorem once in Lean and expose it as a
      reusable API for future OEIS certificates.
\item Build a checker that accepts a proposed defect set, base prefix, and
      periodic forcing table, and returns a machine-checkable recurrence
      certificate.
\item Search OEIS systematically for greedy subset-sum sequences whose
      conjectural recurrences can be resolved by this certificate pattern.
\end{enumerate}

\section{Conclusion}\label{gdf:sec:conclusion}

The two conjectured OEIS recurrences are manifestations of a stronger and
simpler invariant.  In shifted signed-sum coordinates, the reachable set is a
complete interval except for a finite reflected fringe.  Adding the next
greedy term creates three translated copies.  Once the copies overlap without
a high-hole/low-hole collision, the same fringe reproduces itself.  The
arithmetic progression then reads that fringe through a finite residue table,
forcing a periodic correction to doubling and hence a rational recurrence.

For A139217 the invariant has one fringe element, $H=\{3\}$; for A139218 it
has four, $H=\{1,3,6,11\}$.  This proves the complete signed supports, the
periodic deviations, the recurrences, closed forms, generating functions,
asymptotics, and inverses.  It also identifies and repairs the current
A139217 recurrence range; for A139218 it proves the conjecture exactly as
recorded.  The same invariant yields exact counting functions,
cross-sequence identities, and convergent inverse expansions \src{43}.  The
broader finite-fringe method is likely to be more valuable than either
isolated formula, because it reduces an infinite greedy problem to one exact
seed computation and one finite collision check.  More broadly, the results
suggest that congruence-restricted greedy dissociated sets may often be
governed by finite boundary automata; establishing that classification would
convert a wide class of experimentally observed OEIS recurrences into
certifiable theorems \src{43}.

\appendix

\section{Seed support certificates}\label{gdf:app:seeds}

\subsection{A139217}

For $A=\{1,4,7\}$, every integer in $[0,12]$ except $3$ has a ternary-digit
representation:
\begin{center}
\begin{tabular}{c|l@{\qquad}c|l}
\toprule
$x$ & representation & $x$ & representation\\
\midrule
0 & $0$ & 1 & $1$\\
2 & $2\cdot1$ & 4 & $4$\\
5 & $4+1$ & 6 & $4+2\cdot1$\\
7 & $7$ & 8 & $7+1$\\
9 & $7+2\cdot1$ & 10 & $2\cdot4+2\cdot1$\\
11 & $7+4$ & 12 & $7+4+1$\\
\bottomrule
\end{tabular}
\end{center}
The value $3$ is impossible because the contributions from $1$ total at most
$2$, while the next available summand is $4$.  Reflection $x\mapsto24-x$
then proves \eqref{gdf:eq:217-seed-E}.  The signed form
\eqref{gdf:eq:u-base} was proved explicitly in Section~\ref{gdf:sec:A139217}
\src{43}.

\subsection{A139218}

For $A=\{2,5,8\}$, the represented values in $[0,15]$ are certified by
\begin{center}
\begin{tabular}{c|l@{\qquad}c|l}
\toprule
$x$ & representation & $x$ & representation\\
\midrule
0 & $0$ & 2 & $2$\\
4 & $2\cdot2$ & 5 & $5$\\
7 & $5+2$ & 8 & $8$\\
9 & $5+2\cdot2$ & 10 & $2\cdot5$\\
12 & $8+2\cdot2$ & 13 & $8+5$\\
14 & $2\cdot5+2\cdot2$ & 15 & $8+5+2$\\
\bottomrule
\end{tabular}
\end{center}
A direct check of the four remaining values gives the lower fringe
$\{1,3,6,11\}$; reflection $x\mapsto30-x$ gives the four upper holes.
This proves \eqref{gdf:eq:218-seed-E}.

\subsection{Successive support addition \src{43}}

Source 43 summarizes its hand certificate for A139218 by successive support
addition.  Starting from
\[
 D_0=\{0\},\qquad
 D_{k}=D_{k-1}+\{-a_k,0,a_k\},
\]
with $(a_1,a_2,a_3,a_4)=(2,5,8,14)$, one obtains
\begin{align*}
 D_3&=[-15,15]\setminus\{\pm4,\pm9,\pm12,\pm14\},\\
 D_4&=[-29,29]\setminus\{\pm18,\pm23,\pm26,\pm28\}.
\end{align*}
The verification script \file{code/43-boundary-verify.py} evaluates the
same equalities directly.  No unproved numerical assumption enters the
induction.

\section{Derivation of the rational generating functions}\label{gdf:app:gf}

Suppose $a(n)=a(n-1)+a(n-2)+2a(n-3)$ for $n\ge N$.  With
$A(x)=\sum_{n\ge1}a(n)x^n$,
\[
 (1-x-x^2-2x^3)A(x)
\]
automatically has degree at most $N-1$.

For A139217, using $a(1),\dots,a(5)=1,4,7,13,22$ gives residuals
\[
 1,\ 3,\ 2,\ 0,\ -6,
\]
which yields \eqref{gdf:eq:217-ogf}.  For A139218, using
$a(1),\dots,a(6)=2,5,8,14,23,41$ gives residuals
\[
 2,\ 3,\ 1,\ -3,\ -9,\ -12,
\]
which yields \eqref{gdf:eq:218-ogf}.

Conversely, coefficient extraction from either rational function gives the
listed initial values and the common recurrence after the numerator degree.
Thus the generating functions are equivalent to the recurrence plus its
transient data.  (Source 43 argues the same way: multiply each power series
by $1-x-x^2-2x^3$; all coefficients vanish once the recurrence range begins,
and expanding the two rational functions reproduces the OEIS prefixes.)

\section{Audit checklist}\label{gdf:app:checklist}

The logical dependencies of the report are intentionally short.
\begin{enumerate}[label=(\arabic*)]
\item Lemma~\ref{gdf:lem:diss-equivalence} turns greedy admissibility into
      nonmembership in a signed difference set.
\item Lemma~\ref{gdf:lem:E-translation} shifts the signed set to a nonnegative
      ternary support.
\item The exact seed computations identify $H_{217}$ and $H_{218}$.
\item Theorem~\ref{gdf:thm:fringe-propagation} preserves each fringe.
\item Proposition~\ref{gdf:prop:residue-rule} converts the fringe into a finite
      residue table.
\item Theorem~\ref{gdf:thm:periodic-principle} converts periodic corrections into
      the recurrence.
\item Elementary induction gives the closed forms, generating functions, and
      inverse formulas.
\end{enumerate}
No asymptotic assumption enters the proof of either OEIS result.

\section{Tables of terms, partial sums, and defects \src{43}}\label{gdf:app:tables}

\begin{longtable}{r r r r}
\caption{Initial data for A139217.  From $n=3$ onward, the last column is the complete list of positive boundary defects.}\label{gdf:tab:data-u}\\
\toprule
$n$&$u_n$&$U_n$&positive missing value\\
\midrule
\endfirsthead
\toprule
$n$&$u_n$&$U_n$&positive missing value\\
\midrule
\endhead
1&1&1&--\\
2&4&5&--\\
3&7&12&9\\
4&13&25&22\\
5&22&47&44\\
6&49&96&93\\
7&97&193&190\\
8&190&383&380\\
9&385&768&765\\
10&769&1537&1534\\
11&1534&3071&3068\\
12&3073&6144&6141\\
13&6145&12289&12286\\
14&12286&24575&24572\\
\bottomrule
\end{longtable}

\begin{longtable}{r r r l}
\caption{Initial data for A139218.  From $n=3$ onward, the last column is the complete list of positive boundary defects, in increasing order.}\label{gdf:tab:data-v}\\
\toprule
$n$&$v_n$&$V_n$&positive missing values\\
\midrule
\endfirsthead
\toprule
$n$&$v_n$&$V_n$&positive missing values\\
\midrule
\endhead
1&2&2&--\\
2&5&7&--\\
3&8&15&4, 9, 12, 14\\
4&14&29&18, 23, 26, 28\\
5&23&52&41, 46, 49, 51\\
6&41&93&82, 87, 90, 92\\
7&92&185&174, 179, 182, 184\\
8&179&364&353, 358, 361, 363\\
9&353&717&706, 711, 714, 716\\
10&716&1433&1422, 1427, 1430, 1432\\
11&1427&2860&2849, 2854, 2857, 2859\\
12&2849&5709&5698, 5703, 5706, 5708\\
13&5708&11417&11406, 11411, 11414, 11416\\
14&11411&22828&22817, 22822, 22825, 22827\\
\bottomrule
\end{longtable}

Thus the $u$ sequence has the single positive defect $U_n-3$, whereas the $v$
sequence has the four positive defects $V_n-11,V_n-6,V_n-3,V_n-1$.
\mergenote{The intake recomputed both tables from the greedy definition
(brute-force signed difference sets); every entry agrees.}

\section{Provenance and merge record}\label{gdf:app:prov}

\begin{center}
\small
\begin{tabular}{>{\raggedright\arraybackslash}p{1.2cm}>{\raggedright\arraybackslash}p{4.0cm}>{\raggedright\arraybackslash}p{3.3cm}>{\raggedright\arraybackslash}p{1.8cm}>{\raggedright\arraybackslash}p{3.3cm}}
\toprule
Source & Archive (arrival commit) & Main file, PDF & Pin & Author line\\
\midrule
50 (base) & \file{finite_fringe_oeis_A139217_A139218.zip} (\texttt{3a9518c52}) &
  \file{finite_fringe_dissociated_sequences.tex}, 20 pp. &
  \texttt{c79d64038} & ProveIt research report (AI-assisted), prepared for
  Vladimir Reshetnikov with OpenAI\\
43 & \file{ProveIt_OEIS_A139217_A139218.zip} (\texttt{bdf1a1a73}) &
  \file{article.tex}, 24 pp. & none & Research draft prepared for Vladimir
  Reshetnikov and the ProveIt project; computational assistant: OpenAI
  GPT-5.6 Sol Pro\\
\bottomrule
\end{tabular}
\end{center}

Both are manuscripts of batch 73 (cluster O1) of ProveIt's incoming-reports
intake, numbered 50 and 43 there; the placement commit \texttt{9df4ba51a}
staged their files and deleted the archives, which survive in their arrival
commits.  Source 43 was written first (its archive members are timed 17:23 to
17:30 UTC on 1 October 2026; source 50's 18:33 to 18:39).  Source 50's pin
\texttt{c79d64038} already contains source 43's archive, but source 50 never
cites it, and the longest passage the two texts share is 17 words of
\verb|cases| boilerplate or bibliography: they are independent.

\emph{What each contributed.}  Source 50: the base text, in its order, and
everything in it.  Source 43: the second proofs, the threshold form, the
signed certificates, the partial-sum forcing principle, the cross-identities,
the defect counts, the forward and spectral forms, the remainder bound and
explicit inverse coefficients, the counting proofs and asymptotics, the
second verifier, its OEIS wording, roadmap and questions, and the tables.

\emph{Where the merge chose.}
\begin{itemize}
\item Base: source 50, the more general manuscript; its propagation theorem
      contains source 43's.
\item Notation: source 50's throughout; source 43's symbols renamed as in
      Section~\ref{gdf:sec:notation}, including the sign change $d_n\mapsto
      -c_n$, $e_n\mapsto-c_n$.
\item The A139218 counting function uses source 43's residue-zero term, which
      is defined for every real $x$ (merge note in
      Section~\ref{gdf:sec:counting}).
\item Source 43's ``sharp'' wording for the A139217 recurrence range and its
      $n\ge4$ ranges for $V_n$ and $14V_n-13U_n$ are kept with merge notes
      giving the exact ranges.
\item Source 50's sentence on the A290268 report is kept with a merge note.
      Its hard-coded ``Section~6'' became a cross-reference.  Its uncited
      bibliography items (the two OEIS entries and the repository) are now
      cited.
\item Source 43's title page, abstract, executive-summary box, package
      manifest and \verb|listings| styling are not reproduced; their content
      is in Sections~\ref{gdf:sec:intro} and~\ref{gdf:sec:audit} and in the
      README.  Its verbatim code excerpt is printed with \verb|verbatim|.
\item The two bibliographies are merged.
\end{itemize}

\begin{thebibliography}{99}

\bibitem{Bohman1998}
T. Bohman,
\emph{A construction for sets of integers with distinct subset sums},
Electron. J. Combin. \textbf{5} (1998), Research Paper 3.
\url{https://doi.org/10.37236/1341}

\bibitem{CandelaHelfgott2015}
P. Candela and H. A. Helfgott,
\emph{On the dimension of additive sets},
Acta Arith. \textbf{167} (2015), no.~1, 91--100.
\url{https://doi.org/10.4064/aa167-1-5}

\bibitem{CantorMills1966}
D. G. Cantor and W. H. Mills,
\emph{Determination of a subset from certain combinatorial properties},
Canad. J. Math. \textbf{18} (1966), 42--48.
\url{https://doi.org/10.4153/CJM-1966-007-2}

\bibitem{DubroffFoxXu2021}
Q. Dubroff, J. Fox, and M. W. Xu,
\emph{A note on the Erd\H{o}s distinct subset sums problem},
SIAM J. Discrete Math. \textbf{35} (2021), 322--324.
\url{https://doi.org/10.1137/20M1385883}

\bibitem{Dutta2026}
S. Dutta,
\emph{The Greedy Algorithm for Dissociated Sets},
arXiv:2601.07068v4 [math.CO], 5 February 2026.
\url{https://arxiv.org/abs/2601.07068}

\bibitem{Lunnon1988}
W. F. Lunnon,
\emph{Integer sets with distinct subset-sums},
Math. Comp. \textbf{50} (1988), no.~181, 297--320.

\bibitem{OEIS217}
J. W. Layman and The OEIS Foundation,
\emph{A139217: Smallest positive integer of the form $3k+1$ such that all
subsets have different sums}, accessed 1 October 2026.
\url{https://oeis.org/A139217}

\bibitem{OEIS218}
J. W. Layman and The OEIS Foundation,
\emph{A139218: Smallest positive integer of the form $3k+2$ such that all
subsets have different sums}, accessed 1 October 2026.
\url{https://oeis.org/A139218}

\bibitem{ProveIt}
V. Reshetnikov,
\emph{ProveIt}, GitHub repository, inspected at commit
\nolinkurl{c79d640385588fda4b9fdb0cab9fce44790fd100}, 1 October 2026
(source 50; source 43 cites the repository without a commit).
\url{https://github.com/VladimirReshetnikov/ProveIt}

\bibitem{ProveItTransseries}
V. Reshetnikov and contributors,
\emph{Transseries: the polynomial-logarithmic calculus and its inversions},
ProveIt repository, canonical consolidated volume, 2026,
\nolinkurl{Analysis/Transseries/docs/series-and-transseries/Transseries_And_Inversion}.

\bibitem{Steinerberger2023}
S. Steinerberger,
\emph{Some remarks on the Erd\H{o}s distinct subset sums problem},
Int. J. Number Theory \textbf{19} (2023), 1783--1800.
\url{https://doi.org/10.1142/S1793042123500860}

\bibitem{TaoVu2006}
T. Tao and V. H. Vu,
\emph{Additive Combinatorics},
Cambridge Studies in Advanced Mathematics 105,
Cambridge University Press, 2006.

\end{thebibliography}

\end{document}
